% Included by test_loi_giai_mo_rong.tex; compile the parent with XeLaTeX.
\clearpage
% Keep the exercise subsections numbered 10.1 onward without an introductory heading.
\refstepcounter{section}\label{ot45:start}

\subsection{Phần I.A: Năm bài chứng minh cho KDE}

\subsubsection*{Bài 1.1. KDE là một mật độ xác suất hợp lệ}\label{ot45:kde-density}
\textbf{Đề bài.} Với $h>0$, $K\geq0$ và $\int_{\R}K=1$, chứng minh
$\widehat f_h(x)\geq0$ và $\int_{\R}\widehat f_h(x)dx=1$ với mọi mẫu.

\textbf{Điều kiện và công cụ.} Chỉ cần điều kiện về nhân và $h>0$;
không cần giả thiết mật độ thật trơn, cũng không cần dùng tính độc lập.

\buoc{Bước 1. Kiểm tra tính không âm.}
Mỗi số $K((x-X_i)/h)$ không âm, hệ số $1/(nh)>0$, nên tổng có trọng số không âm.

\buoc{Bước 2. Tính diện tích của một bản nhân.}
Giữ $X_i$ cố định, đặt $u=(x-X_i)/h$, $dx=h\,du$.
Khi $x$ chạy từ $-\infty$ đến $\infty$, $u$ cũng chạy trên toàn trục. Do đó
\[
\int_{\R}K\left(\frac{x-X_i}{h}\right)dx=h\int_{\R}K(u)du=h.
\]

\buoc{Bước 3. Cộng diện tích của $n$ số hạng.}
Tổng hữu hạn được đổi thứ tự với tích phân:
\[
\int_{\R}\widehat f_h(x)dx
=\frac1{nh}\sum_{i=1}^n\int_{\R}K\left(\frac{x-X_i}{h}\right)dx
=\frac{nh}{nh}=1.
\]
\textbf{Kết luận.} $\boxed{\widehat f_h\geq0,\quad\int\widehat f_h=1.}$

\textbf{Lỗi dễ mắc.} Quên $dx=h\,du$; hoặc chỉ chứng minh tích phân bằng 1 rồi kết luận là mật độ dù nhân có phần âm.

\subsubsection*{Bài 1.2. Khai triển Bias của KDE}\label{ot45:kde-bias}
\textbf{Đề bài.} Giả sử $f\in C^2(\R)$, $f''$ bị chặn,
$K\in L^1(\R)$, $\int K=1$, $\int uK(u)du=0$ và
$\int u^2|K(u)|du<\infty$. Chứng minh tại $x$ cố định:
\[
\Bias(\widehat f_h(x))=\frac{h^2}{2}\mu_2(K)f''(x)+o(h^2).
\]
Đối xứng là một điều kiện đủ để moment bậc nhất bằng 0, khi moment đó tồn tại.

\textbf{Công cụ.} Tuyến tính kỳ vọng, đổi biến, Taylor dạng phần dư tích phân và hội tụ bị chặn bởi một hàm khả tích.

\buoc{Bước 1. Viết kỳ vọng thành tích phân.}
Các quan sát cùng phân phối nên
\[
\E\widehat f_h(x)=\frac1{nh}\sum_{i=1}^n
\E K\left(\frac{x-X_i}{h}\right)
=\frac{n}{nh}\E K\left(\frac{x-X_1}{h}\right).
\]
Ở bước này chỉ cần cùng phân phối, chưa cần độc lập.
Với biến ngẫu nhiên có mật độ $f$, quy tắc $\E g(X)=\int g(t)f(t)dt$ cho
\begin{align*}
f_h(x)&=\frac1h\E K\left(\frac{x-X_1}{h}\right)
=\frac1h\int_{\R}K\left(\frac{x-t}{h}\right)f(t)dt\\
&=\int_{\R}K(u)f(x-hu)du.
\end{align*}
Ở dòng cuối, $u=(x-t)/h$, $dt=-h\,du$; dấu âm biến mất khi đảo cận.
Chi tiết: $t=x-hu$; khi $t=-\infty$ thì $u=+\infty$,
khi $t=+\infty$ thì $u=-\infty$. Vì vậy
\[
\frac1h\int_{+\infty}^{-\infty}K(u)f(x-hu)(-h)du
=\int_{-\infty}^{+\infty}K(u)f(x-hu)du.
\]

\buoc{Bước 2. Khai triển và kiểm soát phần dư.}
\textbf{Ý tưởng.} Với $u$ cố định và $h$ nhỏ, $x-hu$ gần $x$.
Taylor diễn tả giá trị tại điểm gần đó bằng giá trị, độ dốc và độ cong tại $x$.
Công thức quen thuộc là
$f(x+v)=f(x)+vf'(x)+v^2f''(x)/2+\text{phần dư}$.
Thay $v=-hu$; khi đó $v^2=h^2u^2$, nên dấu của hạng bậc hai là dấu cộng:
\[
f(x-hu)=f(x)-hu f'(x)+\frac{h^2u^2}{2}f''(x)+r_h(u),
\]
với
\[
r_h(u)=h^2u^2\int_0^1(1-s)\{f''(x-shu)-f''(x)\}ds.
\]
\textbf{Phần dư tích phân đến từ đâu?}
Với một độ dịch cố định $v$, đặt $F(s)=f(x+sv)$ trên $0\leq s\leq1$.
Quy tắc hàm hợp cho $F'(s)=vf'(x+sv)$, $F''(s)=v^2f''(x+sv)$.
Tích phân từng phần cho
\[
\int_0^1(1-s)F''(s)ds
=[(1-s)F'(s)]_0^1+\int_0^1F'(s)ds
=-F'(0)+F(1)-F(0).
\]
Chuyển vế, rồi thay lại $F$, được công thức chính xác
\[
f(x+v)=f(x)+vf'(x)+v^2\int_0^1(1-s)f''(x+sv)ds.
\]
Tách $f''(x+sv)$ thành $f''(x)+[f''(x+sv)-f''(x)]$;
phần thứ nhất có tích phân $f''(x)/2$.
Phần thứ hai chính là phần dư ở trên khi thay $v=-hu$.

\textbf{Vì sao phần dư nhỏ sau khi tích phân theo nhân?}
Với từng $u$ cố định, $r_h(u)/h^2\to0$ do $f''$ liên tục tại $x$.
Nếu $M=\|f''\|_\infty$, thì
$|r_h(u)|/h^2\leq Mu^2$.
Thật vậy, hiệu hai giá trị $f''$ có trị tuyệt đối không quá $2M$ và
$\int_0^1(1-s)ds=1/2$, nên chặn là $u^2\cdot2M\cdot(1/2)=Mu^2$.
Vì $Mu^2|K(u)|$ khả tích, định lý hội tụ bị chặn cho
\[
\frac1{h^2}\int K(u)r_h(u)du\longrightarrow0.
\]
Đây là lý do được đưa phần dư qua dấu tích phân.

\buoc{Bước 3. Dùng các moment của nhân.}
Nhân từng hạng Taylor với $K(u)$ rồi tích phân theo $u$.
Các hệ số $f(x),hf'(x),h^2f''(x)/2$ không phụ thuộc $u$:
\begin{align*}
f_h(x)&=f(x)\int K-hf'(x)\int uK
+\frac{h^2}{2}f''(x)\int u^2K+o(h^2)\\
&=f(x)+\frac{h^2}{2}\mu_2(K)f''(x)+o(h^2).
\end{align*}
Ba phép thay ở dòng cuối lần lượt là
$\int K=1$, $\int uK=0$, $\int u^2K=\mu_2(K)$.
Hạng chứa $f'(x)$ biến mất do moment bậc nhất bằng 0, không phải vì $f'(x)=0$.
\textbf{Kết luận.} Trừ $f(x)$ thu được công thức đề bài.
Nếu $f''(x)=0$, hệ số dẫn đầu triệt tiêu và Bias chỉ còn $o(h^2)$.

\textbf{Lỗi dễ mắc.} Viết $O(h^4)$ chỉ từ giả thiết $C^2$;
hoặc dùng khai triển tại một điểm rồi tích phân phần dư mà không có điều kiện kiểm soát.

\subsubsection*{Bài 1.3. Phương sai chính xác và tiệm cận}\label{ot45:kde-variance}
\textbf{Đề bài.} Giả sử $f$ bị chặn và liên tục tại $x$,
$K\in L^1\cap L^2$, $\int K=1$. Chứng minh
\[
\Var(\widehat f_h(x))
=\frac1{nh}\int K(u)^2f(x-hu)du-\frac{f_h(x)^2}{n}
=\frac{f(x)R(K)}{nh}+o((nh)^{-1}).
\]

\textbf{Công cụ.} Các quan sát độc lập; $\Var Z=\E Z^2-(\E Z)^2$;
đổi biến và hội tụ bị chặn.

\buoc{Bước 1. Rút về phương sai của một quan sát.}
Đặt $Z_i=K((x-X_i)/h)$. Do độc lập,
ta dùng $\Var(cY)=c^2\Var(Y)$ và
$\Var(\sum_iZ_i)=\sum_i\Var(Z_i)$.
Các $Z_i$ cùng phân phối nên tổng cuối có $n$ số hạng bằng nhau:
\[
\Var(\widehat f_h(x))=\frac{1}{n^2h^2}\sum_i\Var Z_i
=\frac{1}{nh^2}\{\E Z_1^2-(\E Z_1)^2\}.
\]

\buoc{Bước 2. Tính hai moment bằng đổi biến.}
Đối với moment bậc hai, bình phương $Z_1$ trước khi lấy kỳ vọng:
\begin{align*}
\E Z_1^2
&=\int_{-\infty}^{+\infty}K\left(\frac{x-t}{h}\right)^2f(t)dt\\
&=\int_{+\infty}^{-\infty}K(u)^2f(x-hu)(-h)du
=h\int_{\R}K(u)^2f(x-hu)du.
\end{align*}
Đổi biến tương tự cho moment bậc nhất:
\[
\E Z_1^2=h\int K(u)^2f(x-hu)du,\qquad \E Z_1=hf_h(x).
\]
Thế vào Bước 1 và chia từng số hạng:
\begin{align*}
\Var(\widehat f_h(x))
&=\frac1{nh^2}\left[h\int K(u)^2f(x-hu)du-h^2f_h(x)^2\right]\\
&=\frac1{nh}\int K(u)^2f(x-hu)du-\frac{f_h(x)^2}{n}.
\end{align*}
Lưu ý $\E Z_1^2$ và $(\E Z_1)^2$ là hai đại lượng khác nhau.

\buoc{Bước 3. Xác định số hạng chính và phần dư.}
\textbf{Ý tưởng.} Khi $h\to0$, với mỗi $u$ cố định ta có
$f(x-hu)\to f(x)$. Ta muốn thay mật độ trong tích phân bằng $f(x)$.
Chặn khả tích dưới đây cho phép đưa giới hạn qua dấu tích phân.
Vì $|K(u)^2f(x-hu)|\leq\|f\|_\infty K(u)^2$ khả tích,
\[
\int K(u)^2f(x-hu)du=f(x)R(K)+o(1).
\]
Mặt khác $|f_h(x)|\leq\|f\|_\infty\|K\|_1$, nên
\[
\frac{f_h(x)^2}{n}=O(1/n)=o((nh)^{-1}),
\qquad\frac{1/n}{1/(nh)}=h\to0.
\]
\textbf{Kết luận.} $\boxed{\Var(\widehat f_h(x))=f(x)R(K)/(nh)+o((nh)^{-1}).}$
Nếu $f(x)=0$, kết luận là $o((nh)^{-1})$, không phải tương đương với một số dương.

\textbf{Lỗi dễ mắc.} Bỏ bình phương của $1/h$ khi lấy phương sai;
hoặc bỏ $-(\E Z_1)^2$ mà chưa giải thích bậc của nó.

\subsubsection*{Bài 1.4. Tối ưu AMSE, AMISE và tốc độ}\label{ot45:kde-optimal}
\textbf{Đề bài.} Cho $X_1,\ldots,X_n$ độc lập cùng mật độ $f$ trên $\R$ và
\[
\widehat f_h(x)=\frac1{nh}\sum_{i=1}^nK\left(\frac{x-X_i}{h}\right),\qquad
f_h=\E\widehat f_h,\quad R(g)=\int_{\R}g^2.
\]
Giả sử $K\in L^1\cap L^2$, $\int K=1$, $\int uK(u)du=0$,
$\int u^2|K(u)|du<\infty$ và $\mu_2(K)=\int u^2K(u)du\ne0$.
Mật độ $f$ bị chặn, thuộc $C^2(\R)\cap W^{2,2}(\R)$, $f''$ bị chặn;
$h=h_n\to0$, $nh_n\to\infty$. Tại điểm xét, $f(x)>0$, $f''(x)\ne0$.
\begin{enumerate}[label=(\alph*)]
\item Từ các số hạng dẫn đầu
$\Bias(\widehat f_h(x))=h^2\mu_2(K)f''(x)/2+o(h^2)$ và
$\Var(\widehat f_h(x))=f(x)R(K)/(nh)+o((nh)^{-1})$,
lập và tối ưu AMSE tại $x$.
\item Chứng minh phân rã MISE chính xác; kiểm soát phần dư theo $L^2$ để suy ra
$\AMISE(h)=h^4\mu_2(K)^2R(f'')/4+R(K)/(nh)$.
Giải thích vì sao AMISE không phải MISE chính xác.
\item Với $R(f'')>0$, tìm $h_{\mathrm{AMISE},*}$, chứng minh cực tiểu toàn cục,
tính hằng số của AMISE nhỏ nhất và tốc độ theo $n$.
\end{enumerate}

\textbf{Công cụ.} Phân rã MSE thành bình phương Bias cộng phương sai;
đạo hàm hàm $Ah^4+B/(nh)$ với $A,B>0$.
$W^{2,2}$ nghĩa là $f$ và các đạo hàm yếu đến bậc hai thuộc $L^2$.
\textbf{Lời giải ý (a).}

\buoc{Bước 1. Tối ưu sai số tại một điểm.}
\textbf{Từ Bias sang AMSE:} bình phương hạng Bias chính cho
$[h^2\mu_2(K)f''(x)/2]^2=h^4\mu_2(K)^2[f''(x)]^2/4$,
rồi cộng hạng phương sai chính. Ta tìm $h>0$ làm tổng này nhỏ nhất.
Đặt $A_x=\mu_2(K)^2[f''(x)]^2/4$, $B_x=f(x)R(K)$. Khi đó
\[
\AMSE_x(h)=A_xh^4+\frac{B_x}{nh},\quad
\AMSE_x'(h)=4A_xh^3-\frac{B_x}{nh^2}.
\]
Để thấy rõ đạo hàm, viết $B_x/(nh)=(B_x/n)h^{-1}$.
Khi đó $(A_xh^4)'=4A_xh^3$ và $[(B_x/n)h^{-1}]'=-(B_x/n)h^{-2}$.
Giải phương trình bằng cách nhân hai vế với $h^2>0$:
\[
4A_xh^3=\frac{B_x}{nh^2}
\ \Longleftrightarrow\ 4A_xh^5=\frac{B_x}{n}
\ \Longleftrightarrow\ h^5=\frac{B_x}{4A_xn}.
\]
Lấy căn bậc năm dương và thay $4A_x=\mu_2(K)^2[f''(x)]^2$, được
\[
\boxed{h_{x,*}=\left[\frac{f(x)R(K)}{n\mu_2(K)^2[f''(x)]^2}\right]^{1/5}.}
\]
Đạo hàm bậc hai $12A_xh^2+2B_x/(nh^3)>0$; hàm tiến ra vô cực khi $h\downarrow0$
hoặc $h\to\infty$, nên đây là cực tiểu toàn cục của AMSE.
Đạo hàm bậc hai dương nghĩa là đạo hàm bậc nhất tăng nghiêm ngặt;
trước nghiệm hàm giảm, sau nghiệm hàm tăng. Vì vậy nghiệm là điểm nhỏ nhất duy nhất.

\buoc{Bước 2. Thế nghiệm để xác định tốc độ và hằng số.}
Tại nghiệm, $B_x/(nh_{x,*})=4A_xh_{x,*}^4$. Vì vậy
hai số hạng của AMSE cộng thành $A_xh_{x,*}^4+4A_xh_{x,*}^4=5A_xh_{x,*}^4$.
Để rút gọn, dùng $(ab)^p=a^pb^p$ và $a^pa^q=a^{p+q}$ với các cơ số dương:
\[
5A_x\left(\frac{B_x}{4A_xn}\right)^{4/5}
=5\,4^{-4/5}A_x^{1/5}B_x^{4/5}n^{-4/5}.
\]
Thay $A_x=\mu_2(K)^2[f''(x)]^2/4$ làm hệ số chứa 4 thành
$4^{-4/5}4^{-1/5}=1/4$, nên
\[
\AMSE_x(h_{x,*})=5A_xh_{x,*}^4
=\frac54\{\mu_2(K)^2[f''(x)]^2\}^{1/5}
\{f(x)R(K)\}^{4/5}n^{-4/5}.
\]
Hai số hạng không bằng nhau: phương sai xấp xỉ bằng bốn lần bình phương Bias.

\textbf{Lời giải ý (b).}
\buoc{Bước 3. Phân rã MISE chính xác trước khi dùng xấp xỉ.}
Tại mỗi $x$, viết
$\widehat f_h-f=(\widehat f_h-f_h)+(f_h-f)$.
Hạng chéo có kỳ vọng 0 vì $\E(\widehat f_h-f_h)=0$.
Để thấy rõ, viết $U=\widehat f_h(x)-f_h(x)$ và $b=f_h(x)-f(x)$,
trong đó $b$ không ngẫu nhiên:
\[
\E(U+b)^2=\E U^2+2b\E U+b^2=\Var(\widehat f_h(x))+b^2.
\]
Suy ra
\[
\E(\widehat f_h(x)-f(x))^2
=\Var(\widehat f_h(x))+(f_h(x)-f(x))^2.
\]
Tonelli cho phép đổi kỳ vọng và tích phân của bình phương không âm:
\[
\boxed{\MISE(h)=\E\|\widehat f_h-f\|_2^2
=\int_{\R}\MSE_x(h)dx
=\|f_h-f\|_2^2+\int_{\R}\Var(\widehat f_h(x))dx.}
\]

\buoc{Bước 4. Kiểm soát bình phương chệch trên toàn trục.}
\textbf{Ý tưởng.} Cần xấp xỉ tốt cho toàn bộ hàm, đo bằng chuẩn $L^2$,
để sau đó bình phương và tích phân. Ký hiệu $f''(\,\cdot-a)$ nghĩa là
hàm $x\mapsto f''(x-a)$, tức đường $f''$ được dịch đi một đoạn $a$.
Taylor dạng tích phân và hai moment đầu của $K$ cho
\[
f_h-f=h^2\int u^2K(u)\int_0^1(1-t)f''(\,\cdot-thu)dt\,du.
\]
Trừ $h^2\mu_2(K)f''/2$, chia cho $h^2$ và lấy chuẩn $L^2$:
\begin{align*}
\left\|\frac{f_h-f}{h^2}-\frac{\mu_2(K)}2f''\right\|_2
&\leq\int |u|^2|K(u)|\int_0^1(1-t)
\|f''(\,\cdot-thu)-f''\|_2dt\,du\\
&\longrightarrow0.
\end{align*}
Giới hạn dùng tính liên tục của phép tịnh tiến trong $L^2$;
chuẩn trong tích phân không vượt $2\|f''\|_2$, còn $u^2|K(u)|$ khả tích.
Vì vậy có phần dư $r_h$ với $\|r_h\|_2=o(h^2)$ sao cho
$f_h-f=h^2\mu_2(K)f''/2+r_h$. Khai triển bình phương chuẩn và dùng
Cauchy--Schwarz cho hạng chéo như sau.
Đặt $g=\mu_2(K)f''/2$, khi đó
\[
\|h^2g+r_h\|_2^2=h^4\|g\|_2^2+2h^2\langle g,r_h\rangle+\|r_h\|_2^2.
\]
Vì $|\langle g,r_h\rangle|\leq\|g\|_2\|r_h\|_2$,
hạng chéo có trị tuyệt đối không quá $2h^2\|g\|_2o(h^2)=o(h^4)$;
hạng cuối cũng là $o(h^4)$. Do đó
\[
\|f_h-f\|_2^2=\frac{h^4}{4}\mu_2(K)^2R(f'')+o(h^4).
\]
Đây là điều kiện tích phân thực sự cần; chỉ có khai triển tại từng điểm chưa đủ.

\buoc{Bước 5. Tính phương sai tích phân và nhận diện AMISE.}
Độc lập của mẫu và đổi biến cho công thức chính xác
\[
\Var(\widehat f_h(x))
=\frac1{nh}\int K(u)^2f(x-hu)du-\frac{f_h(x)^2}{n}.
\]
Tích phân hạng đầu dùng Tonelli và $\int f(x-hu)dx=1$, nên
\[
\frac1{nh}\int_{\R}\int_{\R}K(u)^2f(x-hu)du\,dx
=\frac1{nh}\int_{\R}K(u)^2
\underbrace{\left[\int_{\R}f(x-hu)dx\right]}_{=1}du
=\frac{R(K)}{nh}.
\]
Ở tích phân trong, $t=x-hu$ có $dt=dx$ vì $u$ đang được giữ cố định.
Vì vậy
\[
\int\Var(\widehat f_h(x))dx=\frac{R(K)}{nh}-\frac{R(f_h)}n.
\]
Do $\|f_h\|_2\leq\|K\|_1\|f\|_2$, hạng cuối là $O(1/n)$.
Vì $h\to0$, $1/n=o((nh)^{-1})$. Kết hợp các bước trên:
\[
\MISE(h)=\frac{h^4}{4}\mu_2(K)^2R(f'')+\frac{R(K)}{nh}
+o\left(h^4+\frac1{nh}\right).
\]
Ta gọi tổng hai số hạng chính là
\[
\AMISE(h)=\frac{h^4}{4}\mu_2(K)^2R(f'')+\frac{R(K)}{nh}.
\]
Đẳng thức $\int\MSE_x(h)dx=\MISE(h)$ là chính xác;
thay MSE bằng các số hạng tiệm cận chính mới cho AMISE.

\textbf{Lời giải ý (c).}
\buoc{Bước 6. Đạo hàm và kiểm tra cực tiểu toàn cục.}
Đặt $A=\mu_2(K)^2R(f'')/4>0$, $B=R(K)>0$.
Đây chính là bài toán $Ah^4+(B/n)h^{-1}$ đã giải ở ý (a),
với $[f''(x)]^2$ được thay bằng tích phân $R(f'')$.
\[
\AMISE'(h)=4Ah^3-\frac{B}{nh^2}=0
\quad\Longleftrightarrow\quad h^5=\frac{B}{4An}.
\]
Vì $\AMISE''(h)=12Ah^2+2B/(nh^3)>0$ trên $h>0$ và AMISE
tiến ra vô cực tại cả hai đầu miền, nghiệm là cực tiểu toàn cục duy nhất:
\[
\boxed{h_{\mathrm{AMISE},*}
=\left[\frac{R(K)}{n\mu_2(K)^2R(f'')}\right]^{1/5},}
\]
\buoc{Bước 7. Thế nghiệm để tìm hằng số và tốc độ.}
Tại nghiệm, $B/(nh_*)=4Ah_*^4$, nên
\[
Ah_*^4+\frac{B}{nh_*}=5Ah_*^4
=5A\left(\frac{B}{4An}\right)^{4/5}
=5\,4^{-4/5}A^{1/5}B^{4/5}n^{-4/5}.
\]
Thay lại $A,B$ như ở ý (a) cho
\[
\AMISE(h_{\mathrm{AMISE},*})
=\frac54\{\mu_2(K)^2R(f'')\}^{1/5}R(K)^{4/5}n^{-4/5}.
\]
\textbf{Kết luận.} Cả hai bandwidth đều có bậc $n^{-1/5}$;
AMSE/AMISE tối ưu có bậc $n^{-4/5}$. Với khai triển MISE ở trên,
MISE tại bandwidth AMISE cũng đạt $O(n^{-4/5})$.
Bandwidth tối ưu AMISE không được khẳng định là nghiệm chính xác của MISE hữu hạn mẫu.

\textbf{Lỗi dễ mắc.} Dùng $f''(x)^2$ thay $R(f'')$ khi tối ưu toàn trục;
áp dụng công thức điểm khi $f''(x)=0$; hoặc viết tốc độ bandwidth thành tốc độ rủi ro.

\subsubsection*{Bài 1.5. Nhân cấp cao và tốc độ tổng quát}\label{ot45:highorder}
\textbf{Đề bài.} Cho số nguyên $\ell>2$, $K\in L^1\cap L^2$ thỏa
\[
\int K=1,\quad\mu_j(K)=0\ (1\leq j<\ell),\quad
\mu_\ell(K)\ne0,\quad M_\ell:=\int |u|^\ell|K(u)|du<\infty.
\]
Chứng minh Bias bậc $h^\ell$, rồi tìm bậc bandwidth và MISE.
Phần tại điểm dùng $f\in C^\ell$ với $f^{(\ell)}$ bị chặn;
phần MISE dùng $f\in W^{\ell,2}(\R)$.

\textbf{Công cụ.} Taylor, moment triệt tiêu, bất đẳng thức tam giác tích phân trong $L^2$ và tối ưu hàm hai số hạng.

\buoc{Bước 1. Tính Bias tại một điểm.}
\textbf{Đọc trường hợp bậc 4 trước.} Nếu $\ell=4$, Taylor có các hạng
\[
f(x-hu)=f(x)-hu f'(x)+\frac{h^2u^2}{2}f''(x)
-\frac{h^3u^3}{6}f^{(3)}(x)+\frac{h^4u^4}{24}f^{(4)}(x)+r_h(u).
\]
Sau khi nhân với $K$ và tích phân, ba hạng bậc $1,2,3$ mất đi vì
$\mu_1=\mu_2=\mu_3=0$. Trừ $f(x)$, chỉ còn hạng bậc 4 và phần dư.
Đối với bậc tổng quát, cách làm vẫn như vậy; $j!=1\cdot2\cdots j$
và $f^{(j)}$ là đạo hàm lần thứ $j$.
Với $x$ cố định, Taylor đến bậc $\ell$ cho
\[
f(x-hu)=\sum_{j=0}^{\ell}\frac{(-hu)^j}{j!}f^{(j)}(x)+r_h(u).
\]
Ở bậc tổng quát, phần dư có biểu thức
\[
r_h(u)=\frac{(-hu)^\ell}{(\ell-1)!}
\int_0^1(1-t)^{\ell-1}
\{f^{(\ell)}(x-thu)-f^{(\ell)}(x)\}dt.
\]
Nếu $M=\|f^{(\ell)}\|_\infty$, trị tuyệt đối của hiệu trong ngoặc không quá $2M$,
và $\int_0^1(1-t)^{\ell-1}dt=1/\ell$. Do đó
$|r_h(u)|/h^\ell\leq 2M|u|^\ell/\ell!$.
Với mỗi $u$ cố định, tỉ số phần dư tiến về 0 do tính liên tục của $f^{(\ell)}$.
Nhân chặn với $|K(u)|$ cho một hàm khả tích nhờ $M_\ell<\infty$;
hội tụ bị chặn cho $\int K(u)r_h(u)du=o(h^\ell)$ như Bài 1.2.
Vì các moment từ 1 đến $\ell-1$ bằng 0,
ta có thể viết rõ từng hệ số sau tích phân:
\[
f_h(x)-f(x)=\sum_{j=1}^{\ell}
\frac{(-h)^j}{j!}f^{(j)}(x)\underbrace{\int u^jK(u)du}_{\mu_j(K)}
+o(h^\ell).
\]
Trong tổng, chỉ số $j=\ell$ có thể còn lại. Do đó
\[
\boxed{\Bias(\widehat f_h(x))
=\frac{(-1)^\ell\mu_\ell(K)}{\ell!}h^\ell f^{(\ell)}(x)+o(h^\ell)=O(h^\ell).}
\]

\buoc{Bước 2. Chứng minh một chặn tích phân, thay vì chỉ dùng Bias tại điểm.}
\textbf{Ý tưởng.} Chặn độ lớn của toàn bộ hàm Bias trước, rồi bình phương
để chặn sai số tích phân. Khi dịch biến, tích phân trên toàn trục không đổi:
$\int|f^{(\ell)}(x-a)|^2dx=\int|f^{(\ell)}(y)|^2dy$ với $y=x-a$.
Taylor dạng tích phân, sau khi tích phân theo $K$ và triệt tiêu các moment, cho
\[
f_h-f=\frac{(-h)^\ell}{(\ell-1)!}
\int_{\R}u^\ell K(u)\int_0^1(1-t)^{\ell-1}
f^{(\ell)}(\,\cdot-thu)dt\,du.
\]
Phép tịnh tiến không đổi chuẩn $L^2$. Lấy chuẩn, dùng bất đẳng thức tam giác
và $\int_0^1(1-t)^{\ell-1}dt=1/\ell$, được
\[
\|f_h-f\|_2\leq\frac{M_\ell}{\ell!}\|f^{(\ell)}\|_2h^\ell=:C_\ell h^\ell.
\]
Hệ số giai thừa đến từ
$1/[(\ell-1)!]\cdot1/\ell=1/\ell!$.
Bình phương hai vế cho $\|f_h-f\|_2^2\leq C_\ell^2h^{2\ell}$.
Đây là bước bổ sung cần thiết để suy ra MISE trên toàn trục.
Kết hợp phương sai tích phân chính xác, có
\[
\MISE(h)\leq C_\ell^2h^{2\ell}+\frac{R(K)}{nh}.
\]

\buoc{Bước 3. Tối ưu chặn trên.}
Với $C_\ell>0$, đạo hàm vế phải bằng 0 khi
\[
2\ell C_\ell^2h^{2\ell-1}-\frac{R(K)}{nh^2}=0,
\qquad h_*=
\left[\frac{R(K)}{2\ell C_\ell^2 n}\right]^{1/(2\ell+1)}.
\]
Chi tiết giải phương trình: chuyển hạng âm sang vế phải và nhân với $h^2$:
\[
2\ell C_\ell^2h^{2\ell+1}=\frac{R(K)}n,
\qquad h^{2\ell+1}=\frac{R(K)}{2\ell C_\ell^2n}.
\]
Hàm chặn có đạo hàm bậc hai
$2\ell(2\ell-1)C_\ell^2h^{2\ell-2}+2R(K)/(nh^3)>0$,
nên nghiệm dương trên là cực tiểu duy nhất.
Viết $h_*=c\,n^{-1/(2\ell+1)}$ với $c>0$ không phụ thuộc $n$.
Khi đó $h_*^{2\ell}=c^{2\ell}n^{-2\ell/(2\ell+1)}$ và
\[
\frac1{nh_*}=c^{-1}n^{-1+1/(2\ell+1)}
=c^{-1}n^{-2\ell/(2\ell+1)}.
\]
Thế nghiệm vào cả hai hạng:
\[
h_*^{2\ell}\asymp n^{-2\ell/(2\ell+1)},\quad
\frac1{nh_*}\asymp n^{-2\ell/(2\ell+1)}.
\]
\textbf{Kết luận.} $\boxed{h_*\asymp n^{-1/(2\ell+1)},\quad
\MISE(h_*)=O(n^{-2\ell/(2\ell+1)}).}$
Đây là tối ưu của chặn trên vừa chứng minh.
Tính liên tục của phép tịnh tiến trong $L^2$ còn cho khai triển AMISE với
$A_\ell=\mu_\ell(K)^2R(f^{(\ell)})/(\ell!)^2$ thay $C_\ell^2$,
và cùng tốc độ.

\textbf{Ví dụ cụ thể.} Với $\varphi$ là mật độ chuẩn tắc,
$K_4(u)=(3-u^2)\varphi(u)/2$ có
$\mu_1=\mu_2=\mu_3=0$, $\mu_4=(3\cdot3-15)/2=-3$.
Do đó Bias $=-h^4f^{(4)}(x)/8+o(h^4)$ và tốc độ MISE là $n^{-8/9}$.
Các moment chuẩn $\E Z^2=1$, $\E Z^4=3$, $\E Z^6=15$ cho những phép tính này.
Để kiểm tra từng moment: $\int K_4=(3-1)/2=1$;
$\mu_2(K_4)=(3\E Z^2-\E Z^4)/2=(3-3)/2=0$.
Hai moment lẻ bằng 0 do hàm dưới dấu tích phân là hàm lẻ.
Hệ số Bias bậc 4 là $\mu_4/4!=-3/24=-1/8$;
số mũ rủi ro là $2\ell/(2\ell+1)=8/9$ khi $\ell=4$.

\textbf{Lỗi dễ mắc.} Với $\ell>2$, điều kiện $\mu_2=0$ buộc nhân phải có phần âm:
nếu $K\geq0$, $\int u^2K=0$ buộc $K=0$ hầu khắp nơi ngoài $u=0$, trái $\int K=1$.
Vì thế KDE cấp cao không tự động là mật độ không âm.
Nhân đối xứng không thể có moment đầu tiên khác 0 ở bậc lẻ.

\subsection{Phần I.B: Năm bài chứng minh cho ước lượng chiếu}
Trong nhóm này, miền là $[0,1]$, $f\in L^2([0,1])$.
$\phi_1,\ldots,\phi_D$ là các hàm thực trực chuẩn theo tích phân Lebesgue,
$S_D=\operatorname{span}(\phi_1,\ldots,\phi_D)$ và
\[
a_j=\int_0^1f\phi_j,\quad \widehat a_j=\frac1n\sum_i\phi_j(X_i),\quad
f_D=\sum_{j=1}^Da_j\phi_j,\quad \widehat f_D=\sum_{j=1}^D\widehat a_j\phi_j.
\]
Giả sử $\E\phi_j(X)^2<\infty$ khi tính phương sai.
Hai bài 2.2--2.3 phát biểu trên $\R$; bài có miền $[0,1]$ được xem là trường hợp riêng
bằng cách mở rộng các hàm bằng 0 ngoài miền.

\subsubsection*{Bài 2.1. Hệ số mẫu là ước lượng không chệch}\label{ot45:projection-coeff}
\textbf{Đề bài.} Chứng minh $\E\widehat a_j=a_j$, tính phương sai
và giải thích vì sao $\widehat f_D$ vẫn có thể chệch so với $f$.

\textbf{Công cụ.} Kỳ vọng của hàm theo mật độ, tuyến tính kỳ vọng,
phương sai trung bình các quan sát độc lập.

\buoc{Bước 1. Nhận diện hệ số thật là một kỳ vọng.}
\[
\E\phi_j(X_1)=\int_0^1\phi_j(x)f(x)dx=a_j.
\]
Tích phân tuyệt đối hữu hạn do Cauchy--Schwarz:
$\int|\phi_j f|\leq\|\phi_j\|_2\|f\|_2=\|f\|_2$.

\buoc{Bước 2. Lấy kỳ vọng và phương sai của trung bình.}
Đặt $Y_i=\phi_j(X_i)$ với $j$ cố định.
Khi đó $\widehat a_j=(Y_1+\cdots+Y_n)/n$.
Các $Y_i$ độc lập cùng phân phối, nên
\[
\Var(\widehat a_j)=\frac1{n^2}\sum_{i=1}^n\Var(Y_i)
=\frac1n\Var(Y_1)
=\frac1n\{\E Y_1^2-(\E Y_1)^2\}.
\]
Ở đây $\E Y_1=a_j$ và $\E Y_1^2=\int_0^1\phi_j(x)^2f(x)dx$.
Tương tự, tính tuyến tính cho $\E\widehat a_j=(1/n)\sum_i\E Y_i$. Vì vậy
\[
\E\widehat a_j=\frac1n\sum_i a_j=a_j,\qquad
\Var(\widehat a_j)=\frac1n\left(\int_0^1\phi_j(x)^2f(x)dx-a_j^2\right).
\]

\buoc{Bước 3. Phân biệt hai đối tượng được ước lượng.}
\[
\E\widehat f_D(x)=\sum_{j=1}^Da_j\phi_j(x)=f_D(x).
\]
\textbf{Kết luận.} Hệ số mẫu không chệch cho $a_j$;
$\widehat f_D$ không chệch cho $f_D$, nhưng có chệch $f_D-f$ so với mật độ thật.

\textbf{Lỗi dễ mắc.} Suy từ hệ số không chệch thành $\E\widehat f_D=f$ dù đã cắt chuỗi.

\subsubsection*{Bài 2.2. Tối thiểu hóa hàm tương phản}\label{ot45:contrast}
\textbf{Đề bài.} Cho mẫu độc lập $X_1,\ldots,X_n$ cùng mật độ $f\in L^2(\R)$,
không gian $S_D\subset L^2(\R)$ có cơ sở thực trực chuẩn $\phi_1,\ldots,\phi_D$.
Giả sử $\E\phi_j(X)^2<\infty$. Đặt
\[
a_j=\int_{\R}f\phi_j,\quad \widehat a_j=\frac1n\sum_i\phi_j(X_i),\quad
f_D=\sum_ja_j\phi_j,\quad \widehat f_D=\sum_j\widehat a_j\phi_j,
\]
\[
\gamma_n(g)=\|g\|_2^2-\frac2n\sum_i g(X_i),\qquad
V_D=\E\|\widehat f_D-f_D\|_2^2.
\]
\begin{enumerate}[label=(\alph*)]
\item Chứng minh nghiệm duy nhất của $\min_{g\in S_D}\gamma_n(g)$ là $\widehat f_D$
và tìm giá trị tương phản tại nghiệm. Không áp đặt thêm $g\geq0$.
\item Chứng minh $\E\|\widehat f_D\|_2^2=\|f_D\|_2^2+V_D$.
Vì sao chỉ thay $-\|f_D\|_2^2$ bằng $-\|\widehat f_D\|_2^2$ làm xuất hiện một độ lạc quan $V_D$?
\end{enumerate}

\textbf{Công cụ.} Trực chuẩn và hoàn thành bình phương.
Viết $g$ thay vì $h$ để tránh nhầm hàm ứng viên với bandwidth.
\textbf{Lời giải ý (a).}

\buoc{Bước 1. Viết tương phản theo các hệ số.}
\textbf{Vì sao trực chuẩn giúp tính toán?} Trực chuẩn nghĩa là
$\int\phi_j\phi_k=0$ khi $j\ne k$ và bằng 1 khi $j=k$.
Khai triển bình phương tổng trước khi tích phân:
\[
\left\|\sum_jb_j\phi_j\right\|_2^2
=\sum_j\sum_k b_jb_k\int\phi_j(x)\phi_k(x)dx
=\sum_jb_j^2.
\]
Các tích chéo $j\ne k$ mất đi nhờ tích phân, không cần các hệ số độc lập.
Với $g=\sum_{j=1}^D b_j\phi_j$, trực chuẩn cho $\|g\|_2^2=\sum_jb_j^2$;
đồng thời
\[
\frac1n\sum_i g(X_i)=\sum_j b_j\left(\frac1n\sum_i\phi_j(X_i)\right)
=\sum_jb_j\widehat a_j.
\]

\buoc{Bước 2. Hoàn thành bình phương.}
Với hai số $b,a$, ta có $(b-a)^2=b^2-2ba+a^2$,
nên $b^2-2ba=(b-a)^2-a^2$. Áp dụng riêng cho từng $j$:
\[
\gamma_n(g)=\sum_{j=1}^D(b_j^2-2b_j\widehat a_j)
=\sum_{j=1}^D(b_j-\widehat a_j)^2-\sum_{j=1}^D\widehat a_j^2.
\]
Hạng cuối không phụ thuộc $b_j$; tổng bình phương nhỏ nhất bằng 0,
đạt được duy nhất khi mọi $b_j=\widehat a_j$.

\buoc{Bước 3. Ghi nghiệm và giá trị cực tiểu.}
\[
\boxed{\arg\min_{g\in S_D}\gamma_n(g)=\widehat f_D,\qquad
\gamma_n(\widehat f_D)=-\|\widehat f_D\|_2^2.}
\]
\textbf{Kết luận.} Các hệ số tối ưu chính là các trung bình mẫu.
Kiểm tra thêm bằng đạo hàm: $\partial\gamma_n/\partial b_j=2(b_j-\widehat a_j)$,
ma trận Hessian là $2I_D$, xác định dương.
Ở phép kiểm tra này, đạo hàm riêng theo $b_j$ nghĩa là giữ các $b_k$ khác cố định.
Mỗi số hạng chỉ chứa một $b_j$ và có đạo hàm bậc hai bằng 2;
do đó mọi hướng thay đổi hệ số khỏi nghiệm đều làm tương phản tăng.

\textbf{Lời giải ý (b).}
\buoc{Bước 4. Lấy kỳ vọng bình phương chuẩn của ước lượng mẫu.}
Từ định nghĩa phương sai,
$\Var(\widehat a_j)=\E\widehat a_j^2-(\E\widehat a_j)^2$,
nên $\E\widehat a_j^2=\Var(\widehat a_j)+a_j^2$.
Đây là lý do không thể thay kỳ vọng của bình phương bằng bình phương của kỳ vọng.
Do trực chuẩn và $\E\widehat a_j=a_j$,
\[
\E\|\widehat f_D\|_2^2=\sum_j\E\widehat a_j^2
=\sum_j\{a_j^2+\Var(\widehat a_j)\}
=\|f_D\|_2^2+V_D.
\]
Vậy $\E[-\|\widehat f_D\|_2^2]=-\|f_D\|_2^2-V_D$:
tương phản tại mô hình đã khớp mẫu thấp hơn $-\|f_D\|_2^2$ một lượng $V_D$.
Khi mục tiêu tái cấu trúc là $-\|f_D\|_2^2+V_D$, hiệu chỉnh theo kỳ vọng cần
$2V_D$, vì
\[
\E[-\|\widehat f_D\|_2^2+2V_D]=-\|f_D\|_2^2+V_D.
\]
Đây là đồng nhất thức tại $D$ cố định; $V_D$ vẫn chứa phân phối chưa biết.
Bài \hyperref[ot45:oracle]{III.3} giải thích cách chuyển sang mức phạt từ dữ liệu.

\textbf{Lỗi dễ mắc.} Tối ưu trên tập mật độ không âm rồi dùng ngay nghiệm trên toàn $S_D$;
hai bài toán có ràng buộc khác nhau.

\subsubsection*{Bài 2.3. Pythagoras, phân rã và tái cấu trúc rủi ro}\label{ot45:pythagoras}
\textbf{Đề bài.} Cho mật độ $f\in L^2(\R)$, mẫu $X_1,\ldots,X_n$ độc lập cùng mật độ $f$,
và không gian $D$ chiều $S_D\subset L^2(\R)$ có cơ sở thực trực chuẩn
$\phi_1,\ldots,\phi_D$. Dùng tích trong $\langle g,q\rangle=\int_{\R}gq$ và đặt
\[
a_j=\langle f,\phi_j\rangle,\quad \widehat a_j=\frac1n\sum_i\phi_j(X_i),\quad
f_D=P_{S_D}f=\sum_ja_j\phi_j,\quad \widehat f_D=\sum_j\widehat a_j\phi_j.
\]
Giả sử $\E\phi_j(X)^2<\infty$; ký hiệu
$V_D=\E\|\widehat f_D-f_D\|_2^2$ và
$\mathcal R_D=\E\|f-\widehat f_D\|_2^2$.
\begin{enumerate}[label=(\alph*)]
\item Chứng minh đẳng thức Pythagoras với từng mẫu; lấy kỳ vọng để suy ra
phân rã chính xác $\mathcal R_D=\|f-f_D\|_2^2+V_D$. Tính $V_D$ theo các hệ số.
\item Khai triển theo tích trong để chứng minh
$\|f-f_D\|_2^2=\|f\|_2^2-\|f_D\|_2^2$.
Trên một tập ứng viên hữu hạn không rỗng $\mathcal D$, chứng minh
\[
\arg\min_{D\in\mathcal D}\mathcal R_D
=\arg\min_{D\in\mathcal D}\{-\|f_D\|_2^2+V_D\}.
\]
\item Nếu thêm $\sup_{x\in\R}\sum_{j\leq D}\phi_j(x)^2\leq\phi_0D$,
chứng minh $V_D\leq\phi_0D/n$. Phân biệt đẳng thức rủi ro với chặn trên này.
\item Kiểm chứng bằng mật độ $f(x)=2x\mathbf1_{[0,1]}(x)$,
histogram hai ô $[0,1/2)$, $[1/2,1]$ và $n=5$:
tính $\|f\|_2^2$, $\|f_D\|_2^2$, bình phương chệch, $V_D$ và $\mathcal R_D$.
\end{enumerate}

\textbf{Công cụ.} Trực giao của phần dư chiếu, khai triển bình phương chuẩn,
phương sai trung bình mẫu và tính chất bỏ hằng số khỏi bài toán tối ưu.
Miền $[0,1]$ là trường hợp riêng: mở rộng $f$ và các hàm cơ sở bằng 0 ra ngoài miền.

\textbf{Lời giải ý (a).}
\buoc{Bước 1. Chứng minh phần dư chiếu trực giao với $S_D$.}
Trực giao của hai hàm nghĩa là tích trong bằng 0.
Ta kiểm tra trước với từng hàm cơ sở $\phi_k$;
nếu đúng với tất cả các hàm cơ sở thì đúng với mọi tổ hợp tuyến tính của chúng.
Với mọi $k\leq D$,
\[
\langle f-f_D,\phi_k\rangle
=a_k-\sum_{j=1}^Da_j\langle\phi_j,\phi_k\rangle=a_k-a_k=0.
\]
Vì $f_D-\widehat f_D\in S_D$, nên
$\langle f-f_D,f_D-\widehat f_D\rangle=0$ với từng mẫu.
Ví dụ với hai hệ số, $f_D-\widehat f_D=(a_1-\widehat a_1)\phi_1
+(a_2-\widehat a_2)\phi_2$. Tích trong với $f-f_D$ là
$(a_1-\widehat a_1)\cdot0+(a_2-\widehat a_2)\cdot0=0$.

\buoc{Bước 2. Khai triển bình phương chuẩn và loại hạng chéo.}
\begin{align*}
\|f-\widehat f_D\|_2^2
&=\|(f-f_D)+(f_D-\widehat f_D)\|_2^2\\
&=\|f-f_D\|_2^2+2\langle f-f_D,f_D-\widehat f_D\rangle
+\|f_D-\widehat f_D\|_2^2\\
&=\|f-f_D\|_2^2+\|f_D-\widehat f_D\|_2^2.
\end{align*}
Đẳng thức này đúng trước khi lấy kỳ vọng.

\buoc{Bước 3. Lấy kỳ vọng và tính phần phương sai.}
Vì $\E\widehat a_j=a_j$, ta có $\E\widehat f_D=f_D$.
Bình phương Bias theo $L^2$ do đó chính là $\|f-f_D\|_2^2$.
Trực chuẩn cho
\[
\|f_D-\widehat f_D\|_2^2=\sum_{j=1}^D(\widehat a_j-a_j)^2.
\]
Độc lập của các quan sát cho
\[
V_D=\sum_j\Var(\widehat a_j)
=\frac1n\sum_{j=1}^D\left(\int_{\R}\phi_j(x)^2f(x)dx-a_j^2\right).
\]
Lấy kỳ vọng trong Bước 2:
\[
\boxed{\mathcal R_D=\underbrace{\|f-f_D\|_2^2}_{\text{Bias}^2}
+\underbrace{V_D}_{\text{Variance tích phân}}.}
\]
Các hệ số mẫu không cần độc lập; trực chuẩn đã loại hạng chéo khi tích phân.

\phantomsection\label{ot45:risk-reformulation}
\textbf{Lời giải ý (b): tái cấu trúc Bias--Variance.}
\buoc{Bước 4. Tính tích trong giữa mật độ và hình chiếu.}
\[
\langle f,f_D\rangle
=\left\langle f,\sum_ja_j\phi_j\right\rangle
=\sum_ja_j\langle f,\phi_j\rangle
=\sum_ja_j^2=\|f_D\|_2^2.
\]
Tương đương, $f=f_D+(f-f_D)$ và $\langle f-f_D,f_D\rangle=0$.

\buoc{Bước 5. Khai triển bình phương Bias.}
\begin{align*}
\|f-f_D\|_2^2
&=\langle f-f_D,f-f_D\rangle\\
&=\|f\|_2^2-2\langle f,f_D\rangle+\|f_D\|_2^2\\
&=\boxed{\|f\|_2^2-\|f_D\|_2^2}.
\end{align*}
Dấu trừ không làm kết quả âm vì phép chiếu thỏa $\|f_D\|_2\leq\|f\|_2$.

\buoc{Bước 6. Loại hằng số khỏi bài toán chọn $D$.}
Thế vào phân rã rủi ro:
\[
\mathcal R_D=\|f\|_2^2+Q(D),\qquad Q(D)=-\|f_D\|_2^2+V_D.
\]
Với mọi hai ứng viên $D_1,D_2$, ta có
$\mathcal R_{D_1}-\mathcal R_{D_2}=Q(D_1)-Q(D_2)$.
Do đó thứ tự so sánh không đổi, kể cả các trường hợp bằng nhau:
\[
\boxed{\arg\min_{D\in\mathcal D}\mathcal R_D
=\arg\min_{D\in\mathcal D}\{-\|f_D\|_2^2+V_D\}.}
\]
$Q(D)$ có thể âm vì đã trừ hằng số $\|f\|_2^2$; điều đó không có nghĩa rủi ro âm.
Biểu thức vẫn cần $f_D,V_D$ chưa biết. Thay bằng đại lượng mẫu cần hiệu chỉnh
theo Bài \hyperref[ot45:contrast]{2.2} và \hyperref[ot45:oracle]{III.3}.

\textbf{Lời giải ý (c).}
\buoc{Bước 7. Chặn phương sai bằng kích thước mô hình.}
Bỏ các số hạng $-a_j^2\leq0$ và dùng giả thiết cơ sở:
\[
V_D\leq\frac1n\int_{\R}\left(\sum_j\phi_j(x)^2\right)f(x)dx
\leq\frac{\phi_0D}{n}\int_{\R}f(x)dx=\frac{\phi_0D}{n}.
\]
Vậy $\mathcal R_D\leq\|f-f_D\|_2^2+\phi_0D/n$.
Chặn này có thể dùng để tối ưu quy mô $D$; không thay $V_D$ bằng $\phi_0D/n$
trong một đẳng thức hay khẳng định hai bài toán có cùng nghiệm chính xác.

\textbf{Lời giải ý (d).}
\buoc{Bước 8. Tính các đại lượng cho histogram hai ô.}
Với $\phi_j=\sqrt2\mathbf1_{B_j}$,
$p_1=\int_0^{1/2}2x\,dx=1/4$, $p_2=3/4$.
Do đó $f_D$ có chiều cao $2p_1=1/2$, $2p_2=3/2$.
Để thấy chiều cao, $a_j=\int f\sqrt2\mathbf1_{B_j}=\sqrt2p_j$;
trên ô $B_j$ chỉ còn một hàm cơ sở, nên $f_D=a_j\sqrt2=2p_j$.
Các tích phân đầu tiên tính bằng nguyên hàm của đa thức:
\[
p_1=[x^2]_0^{1/2}=\frac14,\qquad
p_2=[x^2]_{1/2}^1=\frac34,\qquad
\int_0^14x^2dx=\left[\frac{4x^3}{3}\right]_0^1=\frac43.
\]
\[
\|f\|_2^2=\int_0^1 4x^2dx=\frac43,\qquad
\|f_D\|_2^2=\frac12\left[\left(\frac12\right)^2+\left(\frac32\right)^2\right]=\frac54.
\]
Tính trực tiếp bình phương chệch kiểm chứng phép tái cấu trúc:
trên ô trái, $(2x-1/2)^2=4x^2-2x+1/4$, nên
\[
\int_0^{1/2}(2x-1/2)^2dx
=\left[\frac43x^3-x^2+\frac x4\right]_0^{1/2}
=\frac16-\frac14+\frac18=\frac1{24}.
\]
Trên ô phải, đặt $y=x-1/2$; cận trở thành $0,1/2$ và
$2x-3/2=2y-1/2$, nên tích phân cũng bằng $1/24$.
Tổng hai ô là
\[
\int_0^{1/2}(2x-1/2)^2dx+\int_{1/2}^1(2x-3/2)^2dx
=\frac1{12}=\frac43-\frac54.
\]
Vì $\sum_j\phi_j^2=2$ trên miền dữ liệu,
\[
V_D=\frac1n\{2-\|f_D\|_2^2\}=\frac15\left(2-\frac54\right)=\frac3{20},
\]
\[
\mathcal R_D=\frac1{12}+\frac3{20}=\frac7{30},\qquad
Q(D)=-\frac54+\frac3{20}=-\frac{11}{10}.
\]
Kiểm tra $\mathcal R_D-\|f\|_2^2=7/30-4/3=-11/10$.

\textbf{Kết luận.} Pythagoras cho phân rã chính xác;
tái cấu trúc cho cùng bài toán tối ưu sau khi bỏ hằng số.
Chặn $V_D\leq\phi_0D/n$ là một bước riêng có thêm giả thiết.

\textbf{Lỗi dễ mắc.} Dùng $\langle f,f_D\rangle=0$ thay vì
$\langle f-f_D,f_D\rangle=0$; hoặc nhầm $V_D$ với phương sai tại một điểm.

\subsubsection*{Bài 2.4. Chọn số chiều tối ưu và hiểu tốc độ minimax}\label{ot45:dimension}
\textbf{Đề bài.} Giả sử cơ sở phù hợp sao cho
$\|f-f_D\|_2^2\leq C_f^2D^{-2\alpha}$ với $\alpha,C_f>0$.
Tối thiểu hóa
\[
\psi(D)=C_f^2D^{-2\alpha}+\frac{\phi_0D}{n},\qquad D\in\{1,2,\ldots\}.
\]
Giải thích điều kiện để gọi tốc độ thu được là tối ưu minimax.

\textbf{Công cụ.} Tối ưu trên biến liên tục rồi chọn số nguyên;
không suy ra sai số xấp xỉ chỉ từ tên của một lớp trơn và một cơ sở bất kỳ.

\buoc{Bước 1. Giải bài toán liên tục.}
Tạm cho số chiều là biến thực $d>0$ để lấy đạo hàm.
Quy tắc $(d^p)'=pd^{p-1}$ với $p=-2\alpha$ cho
$[C_f^2d^{-2\alpha}]'=-2\alpha C_f^2d^{-2\alpha-1}$;
đạo hàm của $\phi_0d/n$ là $\phi_0/n$.
Với $d>0$,
\[
\psi'(d)=-2\alpha C_f^2d^{-2\alpha-1}+\frac{\phi_0}{n}=0
\quad\Longrightarrow\quad
d_*=\left(\frac{2\alpha C_f^2n}{\phi_0}\right)^{1/(2\alpha+1)}.
\]
Các bước giải phương trình là
\[
\frac{\phi_0}{n}=\frac{2\alpha C_f^2}{d^{2\alpha+1}}
\ \Longleftrightarrow\ d^{2\alpha+1}=\frac{2\alpha C_f^2n}{\phi_0},
\]
rồi lấy lũy thừa $1/(2\alpha+1)$ của hai vế dương.
Đạo hàm bậc hai
$2\alpha(2\alpha+1)C_f^2d^{-2\alpha-2}>0$;
hàm giảm trước $d_*$, tăng sau $d_*$.

\buoc{Bước 2. Chọn số chiều nguyên.}
Nếu $d_*<1$, chọn $D_*=1$.
Nếu $d_*\geq1$, tính $\psi$ tại $\lfloor d_*\rfloor$ và $\lceil d_*\rceil$,
chọn giá trị nhỏ hơn; nếu bằng nhau chọn $D$ nhỏ hơn.
Ví dụ $d_*=3.4$ thì chỉ cần so sánh $\psi(3)$ và $\psi(4)$:
hàm đang giảm ở phía trái $3.4$ và tăng ở phía phải, nên các số nguyên xa hơn
không thể tốt hơn cả hai ứng viên này. Không tự động làm tròn theo khoảng cách.
Nếu đề hạn chế $D\leq D_{\max}$, chặn các ứng viên trong $[1,D_{\max}]$ trước khi so sánh.
Với các hằng số cố định và $n$ lớn, trên tập toàn bộ số nguyên dương,
$D_*\asymp n^{1/(2\alpha+1)}$.

\buoc{Bước 3. Thế quy mô nghiệm vào chặn rủi ro.}
Các hằng số không ảnh hưởng số mũ theo $n$.
Quy tắc $(n^a)^b=n^{ab}$ và $n^a/n=n^{a-1}$ cho
\[
\bigl(n^{1/(2\alpha+1)}\bigr)^{-2\alpha}
=n^{-2\alpha/(2\alpha+1)},\qquad
\frac{n^{1/(2\alpha+1)}}n
=n^{[1-(2\alpha+1)]/(2\alpha+1)}.
\]
Vì $1-(2\alpha+1)=-2\alpha$, hai hạng có cùng bậc:
\[
D_*^{-2\alpha}\asymp n^{-2\alpha/(2\alpha+1)},\qquad
\frac{D_*}{n}\asymp n^{-2\alpha/(2\alpha+1)}.
\]
\textbf{Kết luận.} $\boxed{\E\|f-\widehat f_{D_*}\|_2^2
=O(n^{-2\alpha/(2\alpha+1)}).}$
Tại nghiệm liên tục, hạng phương sai của chặn bằng $2\alpha$ lần hạng xấp xỉ.

\textbf{Vì sao cần giả thiết về cơ sở?} Chẳng hạn, với cơ sở đầy đủ và lớp hệ số
$\sum_{j\geq1}j^{2\alpha}a_j^2\leq C_f^2$, Parseval cho
\[
\|f-f_D\|_2^2=\sum_{j>D}a_j^2
\leq(D+1)^{-2\alpha}\sum_{j>D}j^{2\alpha}a_j^2
\leq C_f^2D^{-2\alpha}.
\]
Đây là một điều kiện đủ cụ thể. Với cơ sở lượng giác, lớp Sobolev thích hợp
cần tương thích điều kiện biên/tuần hoàn.

\textbf{Lỗi dễ mắc.} Phép tối ưu trên chỉ chứng minh một cận trên.
Muốn gọi là tối ưu minimax trên lớp $\mathcal F_\alpha$ còn cần định lý cận dưới
\[
\inf_{\widetilde f}\sup_{f\in\mathcal F_\alpha}
\E_f\|\widetilde f-f\|_2^2\geq c n^{-2\alpha/(2\alpha+1)}.
\]
Cận này không được suy ra bằng việc lấy đạo hàm $\psi$.

\subsubsection*{Bài 2.5. Histogram là ước lượng chiếu lên hàm bước}\label{ot45:histogram}
\textbf{Đề bài.} Chia $[0,1]$ thành $D$ ô $B_j$ có độ rộng $1/D$,
các ô nửa kín nửa hở và ô cuối chứa 1. Đặt
$\phi_j(x)=\sqrt D\,\mathbf1_{B_j}(x)$.
Chứng minh cơ sở này trực chuẩn và phép chiếu mẫu chính là histogram.

\textbf{Công cụ.} Tích phân chỉ báo, đếm quan sát trong ô và công thức hệ số mẫu.

\buoc{Bước 1. Kiểm tra trực chuẩn.}
Hàm chỉ báo $\mathbf1_{B_j}$ bằng 1 trong ô $B_j$, bằng 0 ngoài ô.
Do đó tích phân của nó bằng độ dài ô: $\int\mathbf1_{B_j}=|B_j|=1/D$.
Vì chỉ có giá trị 0 hoặc 1, bình phương của hàm chỉ báo vẫn là chính nó.
Các ô rời nhau nên $\int\phi_j\phi_k=0$ khi $j\ne k$.
Với $j=k$, $\int\phi_j^2=D|B_j|=D/D=1$.

\buoc{Bước 2. Tính hệ số và ghép lại.}
Gọi $N_j=\sum_i\mathbf1_{B_j}(X_i)$. Khi đó
\[
\widehat a_j=\frac{\sqrt D\,N_j}{n},\qquad
\widehat f_D(x)=\sum_j\frac{D N_j}{n}\mathbf1_{B_j}(x).
\]
Ở đây $\widehat a_j\phi_j(x)
=(\sqrt D N_j/n)\sqrt D\mathbf1_{B_j}(x)
=(DN_j/n)\mathbf1_{B_j}(x)$.
Nếu $x$ nằm trong ô $B_j$, chỉ hàm chỉ báo của ô đó bằng 1,
các ô khác đều đóng góp 0.
Với $x\in B_j$ và $h=1/D$,
\[
\boxed{\widehat f_D(x)=\frac{DN_j}{n}=\frac{N_j}{nh},}
\]
đúng công thức chiều cao histogram.

\buoc{Bước 3. Kiểm tra bằng mẫu cụ thể.}
Với $X=(0.1,0.2,0.4,0.8)$ và $D=2$, $N_1=3,N_2=1$.
Hệ số là $(3\sqrt2/4,\sqrt2/4)$; chiều cao là $(3/2,1/2)$.
Tổng diện tích $=(3/2)(1/2)+(1/2)(1/2)=1$.

\textbf{Kết luận.} Histogram là trường hợp đặc biệt của phép chiếu.
Ở cơ sở này, $\sum_j\phi_j(x)^2=D$, nên dùng được $\phi_0=1$.

\textbf{Lỗi dễ mắc.} Coi $\widehat a_j$ là chiều cao ô (còn phải nhân $\sqrt D$);
hoặc suy rằng histogram đạt $n^{-4/5}$ khi $f$ trơn bậc hai.
Cơ sở hàm bước thường chỉ có sai số xấp xỉ bình phương bậc $D^{-2}$,
nên tốc độ thông thường vẫn là $n^{-2/3}$.

\subsection{Phần II: Bài tập tính toán số trực tiếp}

\subsubsection*{Bài II.1. Tính hằng số của bốn nhân}\label{ot45:constants}
\textbf{Đề bài.} Tính $\mu_2(K)=\int_{\R}u^2K(u)du$ và $R(K)=\int_{\R}K(u)^2du$
cho các nhân Uniform, Triangular, Epanechnikov và Gaussian.
Ba nhân đầu bằng 0 ngoài $[-1,1]$; Gaussian có giá đỡ toàn trục.

\textbf{Công cụ.} Tính chẵn của các hàm dưới dấu tích phân;
$\int_0^1u^kdu=1/(k+1)$; tích phân Gaussian.

\textbf{Trước khi tính: vì sao có hệ số 2?}
Với hàm chẵn $g$, tức $g(-u)=g(u)$, hai nửa diện tích bằng nhau:
\[
\int_{-1}^1g(u)du=\int_{-1}^0g(u)du+\int_0^1g(u)du
=2\int_0^1g(u)du.
\]
Các nhân trong bài đều chẵn, nên $u^2K(u)$ và $K(u)^2$ cũng chẵn.
Với ba nhân có giá đỡ $[-1,1]$, phần tích phân ngoài khoảng này bằng 0.
Để tính đa thức, dùng
$\int_0^1u^kdu=[u^{k+1}/(k+1)]_0^1=1/(k+1)$ với $k\geq0$.

\buoc{Bước 1. Nhân Uniform $K(u)=\frac12\mathbf1_{\{|u|\leq1\}}$.}
\[
\mu_2(K)=2\int_0^1\frac{u^2}{2}du=\frac13,\qquad
R(K)=\int_{-1}^1\frac14du=\frac12.
\]
Nhân được chuẩn hóa vì $\int_{-1}^1(1/2)du=1$.
Chẳng hạn,
$2\int_0^1u^2/2\,du=[u^3/3]_0^1=1/3$;
$\int_{-1}^1(1/4)du=[u/4]_{-1}^1=1/4-(-1/4)=1/2$.

\buoc{Bước 2. Nhân Triangular $K(u)=(1-|u|)\mathbf1_{\{|u|\leq1\}}$.}
Trên $[0,1]$, $|u|=u$, nên
\[
u^2K(u)=u^2(1-u)=u^2-u^3,\qquad
K(u)^2=(1-u)^2=1-2u+u^2.
\]
Ta chỉ thay $|u|=u$ sau khi đã chuyển về nửa đoạn dương:
\begin{align*}
\mu_2(K)&=2\int_0^1(u^2-u^3)du
=2\left(\frac13-\frac14\right)=\frac16,\\
R(K)&=2\int_0^1(1-2u+u^2)du
=2\left(1-1+\frac13\right)=\frac23.
\end{align*}
Kiểm tra chuẩn hóa: $2\int_0^1(1-u)du=1$.
Viết đầy đủ bước nguyên hàm và thế cận:
\[
\mu_2(K)=2\left[\frac{u^3}{3}-\frac{u^4}{4}\right]_0^1
=2\left(\frac{4-3}{12}\right)=\frac16,
\]
\[
R(K)=2\left[u-u^2+\frac{u^3}{3}\right]_0^1=\frac23.
\]
Nếu tính trên cả $[-1,1]$, phải giữ $|u|$ hoặc tách hai nửa:
\[
\mu_2(K)=\int_{-1}^0u^2(1+u)du+\int_0^1u^2(1-u)du.
\]
Không dùng $u^2(1-u)$ trên nửa âm, vì khi đó $1-|u|=1+u$.

\buoc{Bước 3. Nhân Epanechnikov $K(u)=\frac34(1-u^2)\mathbf1_{\{|u|\leq1\}}$.}
Trước hết,
$u^2K(u)=(3/4)(u^2-u^4)$ và
$K(u)^2=(9/16)(1-2u^2+u^4)$.
Nhân 2 khi chuyển sang $[0,1]$ tạo các hệ số
$2\cdot3/4=3/2$ và $2\cdot9/16=9/8$.
\begin{align*}
\mu_2(K)&=\frac32\int_0^1(u^2-u^4)du
=\frac32\left(\frac13-\frac15\right)=\frac15,\\
R(K)&=\frac98\int_0^1(1-2u^2+u^4)du
=\frac98\left(1-\frac23+\frac15\right)=\frac35.
\end{align*}
Kiểm tra chuẩn hóa: $(3/2)(1-1/3)=1$.
Các nguyên hàm cần dùng là
\[
\mu_2(K)=\frac32\left[\frac{u^3}{3}-\frac{u^5}{5}\right]_0^1,
\qquad
R(K)=\frac98\left[u-\frac{2u^3}{3}+\frac{u^5}{5}\right]_0^1.
\]
Để kiểm tra phân số cuối: $1/3-1/5=2/15$ và
$1-2/3+1/5=(15-10+3)/15=8/15$.

\buoc{Bước 4. Nhân Gaussian $K(u)=(2\pi)^{-1/2}e^{-u^2/2}$.}
Quy tắc hàm hợp cho
\[
K'(u)=\frac1{\sqrt{2\pi}}e^{-u^2/2}\left(-\frac{u^2}{2}\right)'
=-uK(u).
\]
Do đó $u^2K(u)=-uK'(u)$.
Để tích phân từng phần theo $\int a\,db=ab-\int b\,da$,
chọn $a=u$, $da=du$, $db=K'(u)du$, $b=K(u)$ và giữ dấu trừ ở ngoài.
Vì $K'(u)=-uK(u)$, tích phân từng phần và $uK(u)\to0$ ở vô cực cho
\[
\mu_2(K)=-\int_{\R}uK'(u)du
=-[uK(u)]_{-\infty}^{\infty}+\int_{\R}K(u)du=1.
\]
Bình phương nhân làm số mũ đổi từ $-u^2/2$ thành $-u^2$:
\[
R(K)=\frac1{2\pi}\int_{\R}e^{-u^2}du
=\frac{\sqrt\pi}{2\pi}=\frac1{2\sqrt\pi}.
\]
Đây là tích phân trên toàn trục, không phải nguyên hàm:
$\int_{-\infty}^{+\infty}e^{-u^2}du=\sqrt\pi$ là kết quả tích phân Gauss.
Hàm $e^{-u^2}$ không có nguyên hàm bằng các hàm sơ cấp.
Hệ số cuối rút gọn bằng $\pi=\sqrt\pi\sqrt\pi$.

\textbf{Kết luận.}
\begin{center}
\begin{tabular}{lcc}
\toprule Nhân & $\mu_2(K)$ & $R(K)$\\\midrule
Uniform & $1/3$ & $1/2$\\
Triangular & $1/6$ & $2/3$\\
Epanechnikov & $1/5$ & $3/5$\\
Gaussian & $1$ & $1/(2\sqrt\pi)$\\\bottomrule
\end{tabular}
\end{center}
\textbf{Lỗi dễ mắc.} Với $R(K)$ phải bình phương cả hệ số;
với Gaussian không được chỉ tích phân trên $[-1,1]$.

\subsubsection*{Bài II.2. Tính KDE trên mẫu bốn điểm}\label{ot45:calculation}
\textbf{Đề bài.} Cho $X=(1,2,4,5)$, $n=4$, $h=2$, $x_0=3$.
Tính $\widehat f_2(3)$ với nhân Epanechnikov.

\textbf{Công cụ.} Dùng đúng công thức
$\widehat f_2(3)=(1/8)\sum_iK((3-X_i)/2)$.
Giữ phân số đến phép tính cuối.

\buoc{Bước 1. Chuẩn hóa khoảng cách và lập bảng nhân.}
\begin{center}
\begin{tabular}{cccc}
\toprule $X_i$ & $u_i=(3-X_i)/2$ & Thuộc $[-1,1]$? & $K(u_i)$\\\midrule
1 & $1$ & Có & $\frac34(1-1)=0$\\
2 & $1/2$ & Có & $\frac34(1-1/4)=9/16$\\
4 & $-1/2$ & Có & $9/16$\\
5 & $-1$ & Có & $0$\\\bottomrule
\end{tabular}
\end{center}

\buoc{Bước 2. Cộng và chia đúng hệ số.}
\[
\sum_iK(u_i)=\frac{18}{16}=\frac98,\qquad nh=4\cdot2=8,
\]
\[
\boxed{\widehat f_2(3)=\frac18\cdot\frac98=\frac9{64}=0.140625.}
\]
\textbf{Kết luận.} Giá trị mật độ ước lượng tại 3 là $9/64$;
đây không phải xác suất $P(X=3)$.

\textbf{Lỗi dễ mắc.} Chỉ chia cho hai điểm có trọng số dương.
Mẫu vẫn có $n=4$ quan sát; các trọng số bằng 0 không làm thay đổi $n$.

\subsubsection*{Bài II.3. Suy ra quy tắc tham chiếu chuẩn 1.06}\label{ot45:normal-reference}
\textbf{Đề bài.} Cho mẫu độc lập $X_1,\ldots,X_n$ từ mật độ
$f(x)=\sigma^{-1}\varphi((x-\mu)/\sigma)$, $\sigma>0$, $n\geq2$,
với $\varphi(u)=(2\pi)^{-1/2}e^{-u^2/2}$.
KDE $\widehat f_h(x)=(nh)^{-1}\sum_iK((x-X_i)/h)$ dùng $K=\varphi$,
$\mu_2(K)=1$, $R(K)=1/(2\sqrt\pi)$.
\begin{enumerate}[label=(\alph*)]
\item Tính $f''$ và chứng minh $R(f'')=3/(8\sqrt\pi\sigma^5)$.
\item Từ $\AMISE(h)=h^4\mu_2(K)^2R(f'')/4+R(K)/(nh)$,
suy ra công thức chính xác của bandwidth tối ưu, hệ số làm tròn 1.06 và AMISE nhỏ nhất.
\item Khi $\sigma$ chưa biết, viết quy tắc từ độ lệch chuẩn mẫu $s$.
Tính bandwidth cho $n=100$, $s=2$, ghi rõ hệ số chính xác hay làm tròn được sử dụng.
\end{enumerate}

\textbf{Công cụ.} Đạo hàm mật độ chuẩn, đổi biến $z=(x-\mu)/\sigma$ và
tích phân từng phần cho các moment có trọng số $e^{-z^2}$.
\textbf{Lời giải ý (a).}

\buoc{Bước 1. Tính đạo hàm bậc hai.}
\textbf{Ta đang đạo hàm theo biến nào?} Theo $x$; còn $\mu$ và $\sigma$
là các hằng số. Đặt
\[
z=z(x)=\frac{x-\mu}{\sigma},\qquad
\frac{dz}{dx}=\frac1\sigma,\qquad
f(x)=\frac1\sigma\varphi(z(x)).
\]
Viết $z(x)$ để nhắc rằng $z$ cũng thay đổi khi $x$ thay đổi.

\textbf{Hai quy tắc cần dùng.}
Đạo hàm hàm hợp: $\frac{d}{dx}g(z(x))=g'(z(x))z'(x)$
(đạo hàm hàm ngoài rồi nhân đạo hàm hàm trong).
Đạo hàm tích: $(ab)'=a'b+ab'$.

\textbf{1a. Đạo hàm của $\varphi$ theo đối số $z$.}
Vì $(e^{v(z)})'=e^{v(z)}v'(z)$ và $(-z^2/2)'=-z$, nên
\[
\varphi'(z)
=\frac1{\sqrt{2\pi}}e^{-z^2/2}(-z)
=-z\varphi(z).
\]
Ở đây $\varphi'(z)$ là đạo hàm \emph{theo $z$}.
Khi đạo hàm $\varphi(z(x))$ \emph{theo $x$}, phải nhân thêm $z'(x)=1/\sigma$.

\textbf{1b. Đạo hàm lần thứ nhất theo $x$.}
\begin{align*}
f'(x)
&=\frac1\sigma\,\varphi'(z)\,\frac{dz}{dx}\\
&=\frac1\sigma\,[-z\varphi(z)]\,\frac1\sigma
=-\frac{z\varphi(z)}{\sigma^2}.
\end{align*}

\textbf{1c. Đạo hàm lần thứ hai theo $x$.}
Giữ $-1/\sigma^2$ ở ngoài; bên trong là tích của $z(x)$ và $\varphi(z(x))$:
\begin{align*}
f''(x)
&=-\frac1{\sigma^2}\frac{d}{dx}\bigl[z\varphi(z)\bigr]\\
&=-\frac1{\sigma^2}
  \left[\underbrace{\frac1\sigma\varphi(z)}_{z'\,\varphi(z)}
  +\underbrace{z\,\varphi'(z)\frac1\sigma}_{z\,[\varphi(z(x))]'}\right]\\
&=-\frac1{\sigma^2}
  \left[\frac1\sigma\varphi(z)-\frac{z^2}{\sigma}\varphi(z)\right]\\
&=-\frac1{\sigma^3}(1-z^2)\varphi(z).
\end{align*}
Vậy
\[
\boxed{f''(x)=\frac{z^2-1}{\sigma^3}\varphi(z),
\qquad z=\frac{x-\mu}{\sigma}.}
\]
\textbf{Tự kiểm tra dấu:} tại tâm $x=\mu$ thì $z=0$, nên
$f''(\mu)=-\varphi(0)/\sigma^3<0$, phù hợp với đường mật độ cong xuống ở đỉnh.
\textbf{Điểm dễ quên:} mỗi lần đạo hàm hàm hợp theo $x$ đều có thêm
$1/\sigma$; còn $z^2-1$ xuất hiện từ quy tắc đạo hàm tích.

\buoc{Bước 2. Bình phương rồi tích phân.}
Ký hiệu $R(f'')$ nghĩa là lấy \emph{bình phương đạo hàm bậc hai rồi tích phân}:
\[
R(f'')=\int_{\R}[f''(x)]^2\,dx,
\qquad [f''(x)]^2=\frac{(z^2-1)^2}{\sigma^6}\varphi(z)^2.
\]
Đổi biến $x=\mu+\sigma z$ cho $dx=\sigma\,dz$.
Vì vậy $\sigma^{-6}$ nhân với $\sigma$ thành $\sigma^{-5}$.
Ngoài ra,
\[
\varphi(z)^2=\left(\frac1{\sqrt{2\pi}}e^{-z^2/2}\right)^2
=\frac1{2\pi}e^{-z^2},\qquad (z^2-1)^2=z^4-2z^2+1.
\]
Thế vào tích phân:
\[
R(f'')=\frac1{2\pi\sigma^5}\int_{\R}(z^4-2z^2+1)e^{-z^2}dz.
\]
\textbf{Cách tính ba tích phân còn lại.}
Đặt $I_k=\int_{\R}z^{2k}e^{-z^2}dz$. Ta cần $I_0,I_1,I_2$.
Với $I_0$, dùng kết quả tích phân Gauss
$I_0=\int_{\R}e^{-z^2}dz=\sqrt\pi$; không cần tìm nguyên hàm sơ cấp.

Với $k\geq1$, viết $z^{2k}=z^{2k-1}\cdot z$ rồi tích phân từng phần.
Trong công thức $\int a\,db=ab-\int b\,da$, chọn
\[
a=z^{2k-1},\quad db=ze^{-z^2}dz,
\qquad da=(2k-1)z^{2k-2}dz,\quad b=-\tfrac12e^{-z^2}.
\]
Ta tìm được $b$ vì $(e^{-z^2})'=-2ze^{-z^2}$.
Do $z^{2k-1}e^{-z^2}\to0$ khi $z\to\pm\infty$, số hạng biên bằng 0:
\begin{align*}
I_k
&=\left[-\tfrac12z^{2k-1}e^{-z^2}\right]_{-\infty}^{+\infty}
 +\frac{2k-1}{2}\int_{\R}z^{2k-2}e^{-z^2}dz\\
&=\frac{2k-1}{2}I_{k-1}.
\end{align*}
Lần lượt thay $k=1$ và $k=2$:
\[
I_1=\frac12 I_0=\frac{\sqrt\pi}{2},\qquad
I_2=\frac32 I_1=\frac{3\sqrt\pi}{4}.
\]
Cuối cùng, tích phân cần tính là $I_2-2I_1+I_0$, nên
\[
\boxed{R(f'')=\frac{1}{2\pi\sigma^5}
\left(\frac34-1+1\right)\sqrt\pi
=\frac{3}{8\sqrt\pi\sigma^5}.}
\]

\textbf{Lời giải ý (b).}
\buoc{Bước 3. Thế vào băng thông AMISE.}
Với $\mu_2(K)=1$, $R(K)=1/(2\sqrt\pi)$, Bài 1.4 cho
\begin{align*}
h_{\mathrm{AMISE},*}
&=\left[\frac{1/(2\sqrt\pi)}{n\cdot3/(8\sqrt\pi\sigma^5)}\right]^{1/5}\\
&=\left[\frac1{2\sqrt\pi}\cdot\frac{8\sqrt\pi\sigma^5}{3n}\right]^{1/5}
=\left[\frac{4\sigma^5}{3n}\right]^{1/5}\\
&=\left(\frac43\right)^{1/5}\sigma n^{-1/5}
\simeq1.059224\,\sigma n^{-1/5}.
\end{align*}
Đây là bandwidth tối ưu của AMISE dưới mô hình chuẩn.
Ta đã chia một phân số bằng cách nhân với nghịch đảo;
hai thừa số $\sqrt\pi$ triệt tiêu và $8/2=4$.
Vì $\sigma>0$, $(\sigma^5)^{1/5}=\sigma$;
$n^{-1/5}=1/\sqrt[5]{n}$. Hệ số $(4/3)^{1/5}\simeq1.059224$
được làm tròn thành $1.06$.
Tại nghiệm, hạng phương sai gấp bốn lần bình phương chệch, nên
\[
\AMISE(h_*)=\frac54
\left(\frac{3}{8\sqrt\pi\sigma^5}\right)^{1/5}
\left(\frac1{2\sqrt\pi}\right)^{4/5}n^{-4/5}
=\frac{5\,3^{1/5}}{2^{17/5}\sqrt\pi\,\sigma}n^{-4/5}.
\]
Có thể kiểm tra hệ số cuối theo từng cơ số:
$8^{1/5}=2^{3/5}$, $2^{4/5}$ từ thừa số thứ hai và $4=2^2$ ở ngoài
cho $2^{2+3/5+4/5}=2^{17/5}$ ở mẫu;
$(\sqrt\pi)^{1/5}(\sqrt\pi)^{4/5}=\sqrt\pi$;
$(\sigma^5)^{1/5}=\sigma$.

\textbf{Lời giải ý (c).}
\buoc{Bước 4. Thay đại lượng chưa biết bằng ước lượng.}
Lấy $s^2=(n-1)^{-1}\sum_i(X_i-\overline X)^2$ khi $n\geq2$, ta dùng
\[
\boxed{\widehat h_{\mathrm{NR}}\simeq1.06\,s\,n^{-1/5}.}
\]
Ví dụ $n=100$, $s=2$ cho $\widehat h_{\mathrm{NR}}\simeq0.84399$
khi dùng hệ số đã làm tròn $1.06$.
Chi tiết: $100^{-1/5}=10^{-2/5}\simeq0.398107$,
nên $1.06\cdot2\cdot0.398107\simeq0.84399$.

\textbf{Kết luận.} Hệ số 1.06 là phiên bản làm tròn của $(4/3)^{1/5}$.
Thay $\sigma$ bằng $s$ tạo quy tắc từ dữ liệu, không phải đẳng thức với băng thông oracle.
Phiên bản dùng thang đo bền vững $\min(s,\mathrm{IQR}/1.34)$ và phiên bản hệ số $0.9$
được phân biệt thêm ở Mục \ref{sec:silverman}; phép suy ra trên trực tiếp cho hệ số 1.0592.

\textbf{Lỗi dễ mắc.} Quên hệ số $\sigma$ trong $dx$, làm sai lũy thừa $\sigma^{-5}$;
hoặc dùng quy tắc tham chiếu chuẩn như tối ưu cho mọi mật độ nhiều đỉnh.

\subsection{Phần III: Câu hỏi so sánh và lý thuyết có lời giải}

\subsubsection*{Bài III.1. So sánh ba phương pháp ước lượng mật độ}\label{ot45:comparison}
\textbf{Câu hỏi.} So sánh histogram, KDE và phép chiếu về tham số làm trơn,
hình dạng, tốc độ MISE và hạn chế. Tại sao histogram là phép chiếu nhưng thường chậm hơn KDE bậc hai?

\textbf{Điều kiện.} Các tốc độ dưới đây là tốc độ thường đạt được khi chọn tham số đúng bậc,
với điều kiện trơn, nhân, cơ sở và xử lý biên tương ứng.

\buoc{Bước 1. Đối chiếu các đặc điểm.}
\begin{center}
\small
\begin{tabular}{p{.15\textwidth}p{.23\textwidth}p{.24\textwidth}p{.24\textwidth}}
\toprule Tiêu chí & Histogram & KDE & Phép chiếu\\\midrule
Tham số & Độ rộng ô $h$ hoặc số ô $D$ & Băng thông $h$ & Số chiều $D$\\
Tăng chi tiết & Giảm $h$, tăng số ô & Giảm $h$ & Tăng $D$\\
Hình dạng & Hàm nấc & Theo độ trơn của $K$ & Theo hàm cơ sở; có thể là hàm bước hoặc hàm liên tục\\
MISE & $O(n^{-2/3})$ thông thường & $O(n^{-4/5})$ cho nhân bậc hai, $f$ đủ trơn & $O(n^{-2\alpha/(2\alpha+1)})$ nếu xấp xỉ bình phương bậc $D^{-2\alpha}$\\
Hạn chế & Nhạy với gốc chia ô & Chệch biên, nhạy với $h$ & Có thể âm; phụ thuộc cơ sở\\
\bottomrule
\end{tabular}
\end{center}

\buoc{Bước 2. Giải thích khác biệt về tốc độ.}
\textbf{Cân bằng bậc nghĩa là gì?} Ta tìm quy mô $h$ làm hai hạng sai số
có cùng bậc theo $n$. Đây là cách tìm số mũ; để tìm hệ số tối ưu chính xác,
cần đạo hàm như Bài 1.4.
Histogram có bình phương chệch tích phân thường bậc $h^2$, phương sai bậc $1/(nh)$.
Cân bằng $h^2\asymp1/(nh)$ cho $h\asymp n^{-1/3}$, MISE bậc $n^{-2/3}$.
KDE bậc hai có bình phương chệch tích phân bậc $h^4$:
$h^4\asymp1/(nh)$ cho $h\asymp n^{-1/5}$, MISE bậc $n^{-4/5}$.
Viết rõ phép biến đổi:
\begin{align*}
\text{Histogram: }h^2\asymp\frac1{nh}
&\ \Longrightarrow\ h^3\asymp n^{-1}
\ \Longrightarrow\ h\asymp n^{-1/3},\\
\text{KDE: }h^4\asymp\frac1{nh}
&\ \Longrightarrow\ h^5\asymp n^{-1}
\ \Longrightarrow\ h\asymp n^{-1/5}.
\end{align*}
Sau đó bình phương $h$ của histogram cho $n^{-2/3}$;
lấy lũy thừa bậc 4 của $h$ trong KDE cho $n^{-4/5}$.
Với phép chiếu, bậc xấp xỉ phụ thuộc cơ sở; histogram dùng cơ sở hàm bước nên thường
chỉ tương ứng số mũ xấp xỉ $\alpha=1$, kể cả khi $f$ trơn hơn.

\textbf{Kết luận.} Với $\alpha=2$ và cơ sở phù hợp, phép chiếu đạt cùng bậc $n^{-4/5}$
như KDE bậc hai. Không thể gán tốc độ này cho mọi phép chiếu.

\textbf{Lỗi dễ mắc.} Cho rằng phép chiếu luôn liên tục;
hoặc nhầm số mũ: $n^{-4/5}$ giảm nhanh hơn $n^{-2/3}$ khi $n$ tăng.

\subsubsection*{Bài III.2. Chệch biên và hiệu chỉnh phản xạ}\label{ot45:boundary}
\textbf{Câu hỏi.} $X\geq0$, $f(0+)>0$, $f$ trơn một phía tại 0.
Với nhân đối xứng, không âm, tích phân 1, hãy tính chệch KDE tại 0;
xây dựng ước lượng phản xạ, kiểm tra chuẩn hóa và xác định chệch còn lại.

\textbf{Điều kiện và công cụ.} Giả sử $f$ bị chặn, có đạo hàm phải và đạo hàm bậc hai
bị chặn trên nửa trục, $\int u^2K(u)du<\infty$.
Dùng Taylor một phía và chú ý miền tích phân của dữ liệu là $[0,\infty)$.

\buoc{Bước 1. Tính kỳ vọng KDE thông thường tại biên.}
\textbf{Ý tưởng.} Dữ liệu chỉ ở phía dương, nên tại 0 chỉ có một nửa nhân
được tích phân khi $h$ nhỏ. Ta tính cụ thể phần diện tích đó.
Nhờ đối xứng và đổi biến $u=t/h$,
\[
\E\widehat f_h(0)=\frac1h\int_0^\infty K(-t/h)f(t)dt
=\int_0^\infty K(u)f(hu)du.
\]
Chi tiết đổi biến: $t=hu$, $dt=h\,du$, cận $0,+\infty$ giữ nguyên;
$K(-t/h)=K(-u)=K(u)$ do đối xứng và hệ số $h$ triệt tiêu $1/h$.
Khai triển quanh biên từ bên phải:
\[
f(hu)=f(0+)+hu f'(0+)+\rho_h(u),\qquad
|\rho_h(u)|\leq\frac M2h^2u^2,
\]
với $M$ là một chặn của $|f''|$ trên nửa trục.
Nhân với $K(u)$ rồi tích phân: hạng đầu cho
$f(0+)\int_0^\infty K=f(0+)/2$; hạng tuyến tính cho
$h f'(0+)\int_0^\infty uK(u)du$.
Tích phân phần dư được chặn bởi
\[
\left|\int_0^\infty K(u)\rho_h(u)du\right|
\leq\frac M2h^2\int_0^\infty u^2K(u)du=O(h^2).
\]
Đặt $m_{1,+}=\int_0^\infty uK(u)du$. Taylor một phía cho
\[
\E\widehat f_h(0)=\frac12f(0+)+h f'(0+)m_{1,+}+O(h^2),
\]
\[
\boxed{\Bias(\widehat f_h(0))=-\frac12f(0+)+O(h).}
\]
Giá trị $f(0+)$ được hiểu là giới hạn phải, tránh phụ thuộc cách gán giá trị mật độ tại một điểm.
Chệch có số hạng hằng khác 0, nên không phải $O(h)$ hay $O(h^2)$.
Để lấy Bias, phải trừ $f(0+)$ khỏi kỳ vọng:
$f(0+)/2-f(0+)=-f(0+)/2$.

\buoc{Bước 2. Phản xạ từng quan sát qua 0.}
Với $x\geq0$, định nghĩa
\[
\widehat f_h^{\mathrm{ref}}(x)=\frac1{nh}\sum_i
\left[K\left(\frac{x-X_i}{h}\right)+K\left(\frac{x+X_i}{h}\right)\right].
\]
Với $a=X_i/h\geq0$, hai phép đổi biến cho
\[
\int_0^\infty\frac1hK\left(\frac{x-X_i}{h}\right)dx
=\int_{-a}^\infty K(u)du,
\qquad
\int_0^\infty\frac1hK\left(\frac{x+X_i}{h}\right)dx
=\int_a^\infty K(u)du.
\]
Ở phép thứ nhất $x=0$ cho $u=-a$; ở phép thứ hai $x=0$ cho $u=a$.
Tổng diện tích của hai bản nhân trên nửa trục là
\[
\int_{-a}^\infty K(u)du+\int_a^\infty K(u)du=1,
\]
vì đối xứng cho $\int_{-\infty}^{-a}K=\int_a^\infty K$.
Lấy trung bình trên $i$ suy ra
$\int_0^\infty\widehat f_h^{\mathrm{ref}}(x)dx=1$;
tính không âm cũng được bảo toàn.

\buoc{Bước 3. Xác định bậc chệch sau phản xạ.}
Tại 0, hai bản nhân có giá trị bằng nhau:
\[
\E\widehat f_h^{\mathrm{ref}}(0)=2\int_0^\infty K(u)f(hu)du.
\]
Taylor cho
\[
\boxed{\Bias(\widehat f_h^{\mathrm{ref}}(0))
=2h f'(0+)m_{1,+}+O(h^2).}
\]
So với KDE thường, nhân 2 làm hạng đầu của kỳ vọng thành $f(0+)$.
Khi trừ mật độ thật, hạng đó triệt tiêu; hạng tuyến tính được nhân đôi,
và $2O(h^2)$ vẫn là $O(h^2)$.
Vậy phản xạ loại bỏ chệch hằng; thông thường còn $O(h)$.
Nếu $f'(0+)=0$ thì đạt $O(h^2)$ dưới các điều kiện đã nêu.

\buoc{Bước 4. Kiểm tra bằng mật độ mũ và nhân Uniform.}
Với $f(t)=e^{-t}$ trên $t\geq0$, $f(0+)=1$, $f'(0+)=-1$,
$m_{1,+}=\int_0^1u/2\,du=1/4$.
Với nhân Uniform, $K(u)=1/2$ trên $[0,1]$ và bằng 0 khi $u>1$,
nên
\[
\E\widehat f_h(0)=\frac12\int_0^1e^{-hu}du
=\frac12\left[-\frac1h e^{-hu}\right]_0^1
=\frac{1-e^{-h}}{2h}.
\]
Taylor của hàm mũ là $e^{-h}=1-h+h^2/2+O(h^3)$.
Vì vậy $1-e^{-h}=h-h^2/2+O(h^3)$;
chia cho $2h$ cho $1/2-h/4+O(h^2)$.
Các biểu thức chính xác và khai triển tương ứng là
\[
\E\widehat f_h(0)=\frac{1-e^{-h}}{2h}=\frac12-\frac h4+O(h^2),
\]
\[
\E\widehat f_h^{\mathrm{ref}}(0)=\frac{1-e^{-h}}h=1-\frac h2+O(h^2).
\]
Do đó chệch thường tiến về $-1/2$, chệch sau phản xạ bằng $-h/2+O(h^2)$.

\textbf{Kết luận.} Nêu rõ đang xét KDE thường hay đã hiệu chỉnh khi nói về chệch biên.
Để đạt bậc hai mà không cần $f'(0+)=0$, có thể dùng nhân biên triệt tiêu moment bậc nhất;
xem Mục \ref{sec:boundary}.

\textbf{Lỗi dễ mắc.} Dùng $1/(2nh)$ cho công thức phản xạ trên nửa trục:
KDE chuẩn hóa của mẫu mở rộng $2n$ điểm chỉ có một nửa khối lượng ở nửa trục dương;
cần nhân đôi khi dùng làm mật độ trên $[0,\infty)$.

\subsubsection*{Bài III.3. Oracle, tiêu chuẩn phạt và tính thích ứng}\label{ot45:oracle}
\textbf{Đề bài.} Cho mẫu độc lập $X_1,\ldots,X_n$ cùng mật độ $f\in L^2(\R)$, $n\geq2$,
và họ hữu hạn không rỗng các không gian $S_D$ có cơ sở thực trực chuẩn
$\phi_1,\ldots,\phi_D$. Giả sử $\E\phi_j(X)^2<\infty$. Đặt
\[
a_j=\int f\phi_j,\quad \widehat a_j=\frac1n\sum_i\phi_j(X_i),\quad
f_D=\sum_ja_j\phi_j,\quad \widehat f_D=\sum_j\widehat a_j\phi_j,
\]
\[
V_D=\E\|\widehat f_D-f_D\|_2^2,\quad
\mathcal R_D=\E\|f-\widehat f_D\|_2^2,\quad
\gamma_n(g)=\|g\|_2^2-\frac2n\sum_i g(X_i).
\]
\begin{enumerate}[label=(\alph*)]
\item Từ $\mathcal R_D-\|f\|_2^2=-\|f_D\|_2^2+V_D$,
chứng minh $\E[-\|\widehat f_D\|_2^2+2V_D]=\mathcal R_D-\|f\|_2^2$.
Giải thích vì sao cần hiệu chỉnh $2V_D$.
\item Xây dựng một ước lượng không chệch $\widehat V_D$ từ phương sai mẫu của các
$\phi_j(X_i)$. Viết tiêu chuẩn phạt từ dữ liệu và nêu giới hạn của tính không chệch sau khi chọn mô hình.
\item Với bất đẳng thức oracle cho dưới đây, sai số xấp xỉ
$\|f-f_D\|_2^2\leq C_f^2D^{-2\alpha}$, $C_f,\alpha>0$, và một ứng viên
$D_n\asymp n^{1/(2\alpha+1)}$, suy ra tốc độ thích ứng.
Giải thích vì sao quy tắc chọn mô hình không cần nhập $\alpha$.
\end{enumerate}

\textbf{Điều kiện và công cụ.} Xét họ mô hình ứng viên hữu hạn $\mathcal D_n$.
Khi dùng phạt theo chiều, $\phi_0>0$ là hằng số kiểm soát cơ sở qua
$\sup_x\sum_{j\leq D}\phi_j(x)^2\leq\phi_0D$.
Sử dụng Bài 2.2, Bài 2.4 và một định lý chọn mô hình có dạng
\[
\E\|\widehat f_{\widehat D}-f\|_2^2
\leq C\inf_{D\in\mathcal D_n}
\left\{\|f-f_D\|_2^2+\frac{\phi_0D}{n}\right\}+\frac{C'}n.
\]
Các hằng số $C,C'$ không phụ thuộc $n$.
Định lý này là giả thiết được sử dụng trong bài, không phải kết quả cần chứng minh từ đầu.
Nó phải được bảo đảm cho chính quy tắc phạt đang xét; tính không chệch ở ý (b)
không tự chứng minh bất đẳng thức này.

\textbf{Lời giải ý (a).}
\buoc{Bước 1. Xác định phần thông tin không quan sát được.}
Bài 2.4 cho $d_*=(2\alpha C_f^2n/\phi_0)^{1/(2\alpha+1)}$.
Muốn dùng trực tiếp cần biết $\alpha$ và $C_f$, trong khi mẫu không cho sẵn các đại lượng đó.
Đường $f_D$ và sai số $\|f-f_D\|_2^2$ cũng phụ thuộc mật độ thật.
Oracle là chuẩn so sánh lý tưởng giả định biết những thông tin này.
Để chuyển sang dữ liệu, trực chuẩn và tính không chệch hệ số cho
\[
\E\|\widehat f_D\|_2^2=\sum_j\{a_j^2+\Var(\widehat a_j)\}
=\|f_D\|_2^2+V_D.
\]
Vì vậy
\begin{align*}
\E[-\|\widehat f_D\|_2^2+2V_D]
&=-\|f_D\|_2^2-V_D+2V_D\\
&=-\|f_D\|_2^2+V_D=\mathcal R_D-\|f\|_2^2.
\end{align*}
Nếu chỉ cộng $V_D$, kỳ vọng còn $-\|f_D\|_2^2$, thiếu hạng phương sai của rủi ro.
$2V_D$ là hiệu chỉnh theo kỳ vọng tại mô hình cố định, vẫn chưa tính được khi chưa biết $f$.

\textbf{Lời giải ý (b).}
\buoc{Bước 2. Suy ra một tiêu chuẩn chỉ dùng dữ liệu.}
\textbf{Ý tưởng.} Phương sai của trung bình mẫu bằng phương sai một quan sát
chia cho $n$. Vì thế ta ước lượng phương sai từng $\phi_j(X)$ bằng phương sai mẫu,
chia cho $n$, rồi cộng trên $j$.
Theo Bài 2.2, $\gamma_n(\widehat f_D)=-\|\widehat f_D\|_2^2$.
Với từng $j$, phương sai mẫu
\[
s_j^2=\frac1{n-1}\sum_{i=1}^n(\phi_j(X_i)-\widehat a_j)^2
\]
có kỳ vọng $\Var(\phi_j(X))$. Thật vậy,
đặt $Y_i=\phi_j(X_i)$, $m=\E Y_i$, $v=\Var(Y_i)$ và $\overline Y=\widehat a_j$.
Viết $Y_i-\overline Y=(Y_i-m)-(\overline Y-m)$, bình phương rồi cộng:
\begin{align*}
\sum_i(Y_i-\overline Y)^2
&=\sum_i(Y_i-m)^2
-2(\overline Y-m)\sum_i(Y_i-m)+n(\overline Y-m)^2\\
&=\sum_i(Y_i-m)^2-n(\overline Y-m)^2.
\end{align*}
Dòng cuối dùng $\sum_i(Y_i-m)=n(\overline Y-m)$.
Lấy kỳ vọng: hạng đầu bằng $nv$; hạng bị trừ bằng
$n\E(\overline Y-m)^2=n\Var(\overline Y)=n(v/n)=v$.
Tổng còn $(n-1)v$; chia cho $n-1$ cho $\E s_j^2=v$.
Vì $\Var(\widehat a_j)=\Var(\phi_j(X))/n$, ta được
\[
\widehat V_D=\frac1n\sum_{j=1}^Ds_j^2
=\frac1{n(n-1)}\sum_{i=1}^n\sum_{j=1}^D(\phi_j(X_i)-\widehat a_j)^2,
\qquad\E\widehat V_D=V_D.
\]
Đặt $\widehat Q(D)=-\|\widehat f_D\|_2^2+2\widehat V_D$. Tuy hai hạng dùng cùng mẫu,
tuyến tính kỳ vọng vẫn cho $\E\widehat Q(D)=\mathcal R_D-\|f\|_2^2$ tại từng $D$ cố định.
Đó là một cách xây dựng mức phạt tính được từ dữ liệu.
Viết tổng quát với $\operatorname{pen}(D)$:
\[
\boxed{\widehat D\in\arg\min_{D\in\mathcal D_n}
\left\{-\sum_{j=1}^D\widehat a_j^2+\operatorname{pen}(D)\right\}.}
\]
Một dạng phạt thường xét là $\operatorname{pen}(D)=\kappa\phi_0D/n$.
Chọn $D$ nhỏ nhất nếu bằng điểm. $\kappa$ phải được hiệu chỉnh theo định lý hoặc được đề bài cho.
Trong mô hình lồng nhau, phần âm giảm khi thêm các hệ số;
mức phạt giúp kiểm soát xu hướng chọn mô hình quá lớn.
Phạt $2\widehat V_D$ và phạt $\kappa\phi_0D/n$ là hai cách khác nhau;
không đồng nhất chúng như một công thức chính xác.
Sau khi lấy $\widehat D$ tối thiểu hóa điểm số,
$\widehat Q(\widehat D)$ không tự động là ước lượng không chệch cho rủi ro mô hình đã chọn.
Đánh giá chất lượng lựa chọn cần kiểm soát ngẫu nhiên đồng thời trên họ mô hình.

\textbf{Kiểm chứng với ví dụ histogram hai ô.} Với $f(x)=2x\mathbf1_{[0,1]}(x)$,
$n=5$, đặt $N$ là số quan sát trong ô trái; $N\sim\operatorname{Binomial}(5,1/4)$.
Khi $N=k$, chuẩn bình phương mẫu và ước lượng phương sai là
\[
T_k=\|\widehat f_D\|_2^2=2\left[\left(\frac k5\right)^2+
\left(\frac{5-k}{5}\right)^2\right],\qquad
\widehat V_D=\frac{2-T_k}{4}.
\]
Công thức sau được suy ra từ $\sum_i\sum_j\phi_j(X_i)^2=2n$ và
$\sum_i\sum_j(\phi_j(X_i)-\widehat a_j)^2=n(2-\|\widehat f_D\|_2^2)$.
Để kiểm tra đẳng thức này, khai triển từng bình phương rồi cộng theo $i$:
\[
\sum_i(\phi_j(X_i)-\widehat a_j)^2
=\sum_i\phi_j(X_i)^2-2\widehat a_j\underbrace{\sum_i\phi_j(X_i)}_{n\widehat a_j}
+n\widehat a_j^2
=\sum_i\phi_j(X_i)^2-n\widehat a_j^2.
\]
Cộng theo $j$ cho $2n-nT_k$; chia cho $n(n-1)$ được
$(2-T_k)/(n-1)=(2-T_k)/4$.
Nếu muốn kiểm tra bằng sáu trường hợp, dùng bảng sau;
các xác suất có mẫu số chung $4^5=1024$, nên hàng cuối ghi $1024P(N=k)$:
\[
\begin{array}{c|rrrrrr}
k&0&1&2&3&4&5\\\hline
25T_k&50&34&26&26&34&50\\
1024P(N=k)&243&405&270&90&15&1
\end{array}
\]
Ở mỗi cột, điểm số là $-T_k+2(2-T_k)/4=1-3T_k/2$.
Kỳ vọng hữu hạn cho
\[
\sum_{k=0}^5\binom5k\left(\frac14\right)^k\left(\frac34\right)^{5-k}
\left[-T_k+2\frac{2-T_k}{4}\right]
=-\frac{11}{10}=\mathcal R_D-\|f\|_2^2.
\]

\textbf{Lời giải ý (c).}
\buoc{Bước 3. Dùng bất đẳng thức oracle để chứng minh tốc độ.}
Giả sử trên lớp mật độ xét có chặn đều
$\|f-f_D\|_2^2\leq C_f^2D^{-2\alpha}$,
và họ ứng viên chứa một $D_n\asymp n^{1/(2\alpha+1)}$.
Vì infimum không lớn hơn giá trị tại $D_n$,
ta có thể thay bài toán tìm giá trị nhỏ nhất bằng giá trị của riêng ứng viên này
để lấy một chặn trên. Với $D_n\asymp n^{1/(2\alpha+1)}$,
\[
D_n^{-2\alpha}\asymp n^{-2\alpha/(2\alpha+1)},\qquad
\frac{D_n}{n}\asymp n^{1/(2\alpha+1)-1}
=n^{-2\alpha/(2\alpha+1)}.
\]
Thế hai bậc đó vào bất đẳng thức oracle:
\begin{align*}
\E\|\widehat f_{\widehat D}-f\|_2^2
&\leq C\left(C_f^2D_n^{-2\alpha}+\frac{\phi_0D_n}{n}\right)+\frac{C'}n\\
&=O(n^{-2\alpha/(2\alpha+1)})+O(n^{-1})\\
&=O(n^{-2\alpha/(2\alpha+1)}).
\end{align*}
Dòng cuối dùng $2\alpha/(2\alpha+1)<1$.
Đặt $p=2\alpha/(2\alpha+1)$; khi đó $n^{-1}/n^{-p}=n^{p-1}\to0$.
Vì vậy số hạng $C'/n$ nhỏ hơn về bậc và được gộp vào $O(n^{-p})$.
Quy tắc tính $\widehat D$ không nhập $\alpha$, nhưng vẫn đạt bậc ứng với độ trơn thật:
đó là tính thích ứng trong các lớp mà định lý bảo đảm.

\buoc{Bước 4. Minh họa phép tính điểm số.}
Giả sử $n=100$, $\phi_0=1$, mức phạt đề cho $\kappa=2$,
ứng viên $D\in\{2,4,8\}$ và các chuẩn bình phương mẫu lần lượt là $1.10,1.17,1.20$.
Điểm số là
\[
\begin{array}{c|ccc}
D&2&4&8\\\hline
-\|\widehat f_D\|_2^2+2D/100&-1.06&-1.09&-1.04
\end{array}
\]
Chọn $\widehat D=4$ vì $-1.09$ nhỏ nhất.
Ví dụ tính số này không khẳng định $\kappa=2$ đủ cho mọi định lý oracle.

\textbf{Kết luận.} Tính thích ứng suy ra từ một bảo đảm oracle đã được chứng minh
cho đúng họ mô hình và đúng mức phạt.
Thảo luận về hiệu chỉnh phạt và các bất đẳng thức oracle xem
\href{https://arxiv.org/abs/0910.1654}{Lerasle, \emph{Optimal model selection in density estimation}}.

\textbf{Lỗi dễ mắc.} Tuyên bố mọi $\kappa>0$ đều cho tính thích ứng;
hoặc coi phân rã Pythagoras đã chứng minh bất đẳng thức oracle cho $\widehat D$ phụ thuộc dữ liệu.

\subsubsection*{Bài III.4. LOOCV cho KDE: suy ra tiêu chuẩn}\label{ot45:loocv}
\textbf{Đề bài.} Cho mẫu $X_1,\ldots,X_n$ độc lập cùng mật độ
$f\in L^2(\R)$, $n\geq2$. Nhân $K\in L^1\cap L^2$ đối xứng, $\int K=1$.
Với $h>0$ cố định, đặt
\[
\widehat f_h(x)=\frac1{nh}\sum_iK\left(\frac{x-X_i}{h}\right),\quad
\widehat f_{h,-i}(x)=\frac1{(n-1)h}\sum_{j\ne i}K\left(\frac{x-X_j}{h}\right).
\]
Giữ nguyên $h$ khi bỏ từng quan sát, gọi $f_h=\E\widehat f_h$ và $R(f)=\|f\|_2^2$.
Xét điểm số $\widehat J(h)=\int\widehat f_h^2-(2/n)\sum_i\widehat f_{h,-i}(X_i)$.
\begin{enumerate}[label=(\alph*)]
\item Khai triển rủi ro $\MISE(h)=\E\|\widehat f_h-f\|_2^2$ và nhận diện hạng không phụ thuộc $h$.
\item Tính kỳ vọng có điều kiện của $\widehat f_{h,-i}(X_i)$ theo mẫu còn lại.
Chứng minh $\E\widehat f_{h,-i}(X_i)=\E\int\widehat f_h f$;
biện luận việc đổi tích phân và kỳ vọng.
\item Chứng minh $\E\widehat J(h)=\MISE(h)-R(f)$.
Giải thích đúng ý nghĩa không chệch và giới hạn khi chọn $\widehat h$ từ dữ liệu.
\item Viết tiêu chuẩn dưới dạng tổng đôi, chỉ rõ tổng nào chứa đường chéo $i=j$.
\end{enumerate}

\textbf{Công cụ.} Khai triển bình phương chuẩn, độc lập giữa điểm bỏ ra và mẫu còn lại,
kỳ vọng có điều kiện, Fubini và Cauchy--Schwarz.
Trong lời giải, $\mathrm{CV}(h)$ và $\widehat J(h)$ chỉ cùng một điểm số.

\textbf{Lời giải ý (a).}
\buoc{Bước 1. Khai triển sai số rồi lấy kỳ vọng.}
\[
\|\widehat f_h-f\|_2^2
=\int\widehat f_h^2-2\int\widehat f_h f+R(f).
\]
Lấy kỳ vọng hai vế:
\[
\MISE(h)=\E\int\widehat f_h^2-2\E\int\widehat f_h f+R(f).
\]
Hạng $R(f)$ không đổi theo $h$. Hạng đầu tính được từ mẫu;
hạng chứa $f$ được ước lượng bằng trung bình đánh giá trên các điểm bỏ ra.

\textbf{Lời giải ý (b).}
\buoc{Bước 2. Chứng minh trung bình leave-one-out có đúng kỳ vọng.}
\textbf{Ý tưởng bằng lời.} Tạm giữ nguyên $n-1$ điểm dùng dựng đường KDE,
chỉ cho điểm bị bỏ ra $X_i$ thay đổi. Điểm đó độc lập với đường đã dựng,
nên trung bình giá trị đường tại $X_i$ là tích phân đường nhân mật độ $f$.
Sau đó mới lấy trung bình tiếp trên $n-1$ điểm đã giữ cố định.
Gọi $\mathcal F_{-i}=\sigma(X_j:j\ne i)$.
Ở đây $\mathcal F_{-i}$ ký hiệu toàn bộ thông tin của mẫu còn lại;
ký hiệu $\sigma(\cdot)$ trong dòng này chỉ một tập thông tin, không phải độ lệch chuẩn.
Khi điều kiện theo $\mathcal F_{-i}$, hàm $\widehat f_{h,-i}$ đã cố định,
còn $X_i$ vẫn có mật độ $f$ do độc lập. Vì vậy
\[
\E[\widehat f_{h,-i}(X_i)\mid\mathcal F_{-i}]
=\int\widehat f_{h,-i}(x)f(x)dx.
\]
Với mỗi $x$, tính tuyến tính và các quan sát cùng phân phối cho
\[
\E\widehat f_{h,-i}(x)
=\frac{n-1}{(n-1)h}\E K\left(\frac{x-X_1}{h}\right)
=\int K(u)f(x-hu)du=f_h(x).
\]
Chỉ số 1 ở đây đại diện một quan sát có cùng phân phối, không yêu cầu nó nằm trong mẫu còn lại.
Mẫu kích thước $n-1$ và mẫu kích thước $n$ có cùng kỳ vọng KDE vì đã chuẩn hóa đúng.

\buoc{Bước 3. Biện luận khả tích và lấy kỳ vọng lần nữa.}
\textbf{Mục đích của bước kỹ thuật.} Ta sắp đổi thứ tự kỳ vọng và tích phân.
Để làm vậy ngay cả khi $K$ có phần âm, cần chứng minh tích phân của trị tuyệt đối
có kỳ vọng hữu hạn. Nếu chỉ muốn theo dõi cách xây dựng CV,
có thể đọc kết quả của chặn dưới đây rồi sang Bước 4.
Đặt $K_h(t)=K(t/h)/h$. Do $f\in L^2$,
ký hiệu tích chập $(|K_h|*f)(x)$ nghĩa là
$\int |K_h(x-t)|f(t)dt$.
Bất đẳng thức Cauchy--Schwarz chặn tích phân tích hai hàm bằng tích hai chuẩn;
bất đẳng thức Young cho $\||K_h|*f\|_2\leq\|K_h\|_1\|f\|_2$
và đổi biến cho $\|K_h\|_1=\|K\|_1$. Cụ thể,
Young và Cauchy--Schwarz cho
\[
\int f(x)(|K_h|*f)(x)dx
\leq\|f\|_2\,\||K_h|*f\|_2
\leq\|K\|_1\|f\|_2^2<\infty.
\]
Mặt khác $\E|\widehat f_{h,-i}(x)|\leq(|K_h|*f)(x)$,
nên $\E\int|\widehat f_{h,-i}(x)|f(x)dx<\infty$.
Ước lượng đủ mẫu có cùng chặn. Fubini áp dụng được kể cả khi $K$ có dấu.
Tính lặp của kỳ vọng có điều kiện dẫn đến
\begin{align*}
\E\widehat f_{h,-i}(X_i)
&=\E\int\widehat f_{h,-i}(x)f(x)dx\\
&=\int\E\widehat f_{h,-i}(x)f(x)dx\\
&=\int f_h(x)f(x)dx
=\E\int\widehat f_h(x)f(x)dx.
\end{align*}
Hạng bình phương cũng khả tích: do bất đẳng thức bình phương trung bình,
\[
\|\widehat f_h\|_2^2
\leq\frac1n\sum_i\|K_h(\,\cdot-X_i)\|_2^2=\frac{R(K)}h<\infty.
\]

\textbf{Lời giải ý (c).}
\buoc{Bước 4. Thế đẳng thức kỳ vọng vào điểm số.}
\[
\boxed{\mathrm{CV}(h)=\int\widehat f_h(x)^2dx
-\frac2n\sum_i\widehat f_{h,-i}(X_i),}
\]
\begin{align*}
\E\widehat J(h)
&=\E\int\widehat f_h^2-\frac2n\sum_i\E\widehat f_{h,-i}(X_i)\\
&=\E\int\widehat f_h^2-2\E\int\widehat f_h f\\
&=\boxed{\E\|\widehat f_h-f\|_2^2-\|f\|_2^2}.
\end{align*}
Vì vậy có thể chọn $\widehat h$ tối thiểu hóa $\mathrm{CV}(h)$ trên một lưới bandwidth dương.
Điểm số không chệch cho $\MISE(h)-R(f)$, chứ không cho chính MISE khi chưa cộng lại $R(f)$.
Điểm số có thể âm do đã bỏ hằng số này.
Tính không chệch này đúng với từng $h$ cố định;
không kết luận giá trị nhỏ nhất sau khi chọn $\widehat h$ vẫn là ước lượng không chệch của rủi ro được chọn.

\textbf{Lời giải ý (d).}
\buoc{Bước 5. Biến đổi thành tổng đôi tính được.}
Đặt $C_K(v)=\int K(u)K(u+v)du$ là tự tương quan của nhân.
\textbf{Từ bình phương một tổng sang tổng đôi.}
Quy tắc $(\sum_i a_i)^2=\sum_i\sum_j a_ia_j$ cho
\[
\int\widehat f_h(x)^2dx
=\frac1{n^2h^2}\sum_i\sum_j
\int K\left(\frac{x-X_i}{h}\right)K\left(\frac{x-X_j}{h}\right)dx.
\]
Với từng cặp $(i,j)$, đặt $u=(x-X_i)/h$, $x=X_i+hu$, $dx=h\,du$.
Đối số của nhân thứ hai trở thành
\[
\frac{x-X_j}{h}=u+\frac{X_i-X_j}{h}.
\]
Cận vẫn là toàn trục. Tích phân của tích hai nhân bằng
$hC_K((X_i-X_j)/h)$; nhân với $1/(n^2h^2)$ còn $1/(n^2h)$.
Vì vậy
\[
\int\widehat f_h^2=\frac1{n^2h}\sum_{i,j}C_K\left(\frac{X_i-X_j}{h}\right).
\]
Với nhân đối xứng, $C_K=K*K$.
Thật vậy, đổi biến $t=-u$ cho
$C_K(v)=\int K(-t)K(v-t)dt=\int K(t)K(v-t)dt$.
Thế trực tiếp định nghĩa leave-one-out vào hạng còn lại:
\[
\boxed{\mathrm{CV}(h)=\frac1{n^2h}\sum_{i,j}C_K\left(\frac{X_i-X_j}{h}\right)
-\frac{2}{n(n-1)h}\sum_{i\ne j}K\left(\frac{X_i-X_j}{h}\right).}
\]
Tổng đầu có đường chéo $i=j$, tổng thứ hai loại đường chéo;
mọi tổng đều tính theo cặp có thứ tự.
Ví dụ $n=3$, tổng thứ hai gồm sáu cặp
$(1,2),(1,3),(2,1),(2,3),(3,1),(3,2)$.
Tổng đầu gồm cả sáu cặp này và ba cặp $(1,1),(2,2),(3,3)$.
Với $i=j$, $C_K(0)=\int K^2=R(K)$, nên riêng phần đường chéo
của hạng đầu bằng $nR(K)/(n^2h)=R(K)/(nh)$.

\textbf{Kết luận.} Tiêu chuẩn không cần biết $f$.
Với Gaussian, $C_K(v)=(2\sqrt\pi)^{-1}e^{-v^2/4}$, tức mật độ chuẩn phương sai 2,
nên cả hai tổng đều tính trực tiếp được.
Để thấy công thức Gaussian của $C_K$, hoàn thành bình phương trong số mũ:
\[
u^2+(u+v)^2=2(u+v/2)^2+v^2/2.
\]
Vì vậy
\[
C_K(v)=\frac{e^{-v^2/4}}{2\pi}
\int_{\R}e^{-(u+v/2)^2}du
=\frac{e^{-v^2/4}}{2\pi}\sqrt\pi
=\frac1{2\sqrt\pi}e^{-v^2/4}.
\]
Tích phân ở giữa dùng $w=u+v/2$, $dw=du$ và tích phân Gauss trên toàn trục.

\textbf{Lỗi dễ mắc.} Dùng $n$ thay $n-1$ trong ước lượng bỏ một điểm;
hoặc chỉ giữ hạng tích chập và hiệu chỉnh $K(0)$ mà làm mất tổng đôi thứ hai.
Mục \ref{sec:cv} có thêm ví dụ tính tay LOOCV bằng phân số.
